\documentclass[11pt]{article}

\usepackage[margin=1in]{geometry}
\usepackage{graphicx}
\usepackage{booktabs}
\usepackage{multirow}
\usepackage{amsmath}
\usepackage{amssymb}
\usepackage{xcolor}
\usepackage{caption}
\usepackage{subcaption}
\usepackage[hidelinks]{hyperref}

% Figures live next to the checkpoints; point graphicspath at the repo root.
\graphicspath{{../}}

% Short handles for each run's figure folder.
\newcommand{\runA}{checkpoints/dualvae/dualvae_20260724-121508_392b5b}
\newcommand{\runB}{checkpoints/dualvae/dualvae_20260724-133147_b03f09}
\newcommand{\runC}{checkpoints/dualvae/dualvae_20260725-015430_5878cc}
\newcommand{\runD}{checkpoints/dualvae/dualvae_20260725-134828_ce6208}
\newcommand{\tagA}{dualvae_20260724-121508_392b5b}
\newcommand{\tagB}{dualvae_20260724-133147_b03f09}
\newcommand{\tagC}{dualvae_20260725-015430_5878cc}
\newcommand{\tagD}{dualvae_20260725-134828_ce6208}
\newcommand{\runAg}{checkpoints/dualvae/dualvae_20260801-201243_0b2ee6}
\newcommand{\runBg}{checkpoints/dualvae/dualvae_20260801-201309_1e27ab}
\newcommand{\tagAg}{dualvae_20260801-201243_0b2ee6}
\newcommand{\tagBg}{dualvae_20260801-201309_1e27ab}
\newcommand{\runCg}{checkpoints/dualvae/dualvae_20260726-230719_5e588d}
\newcommand{\runDg}{checkpoints/dualvae/dualvae_20260726-230704_866e33}
\newcommand{\tagCg}{dualvae_20260726-230719_5e588d}
\newcommand{\tagDg}{dualvae_20260726-230704_866e33}
% E+ -- standard prior pushed one decade past B+ (beta = 1.0), same LPIPS/GAN recipe.
\newcommand{\runEg}{checkpoints/dualvae/dualvae_20260802-163932_ca522b}
\newcommand{\tagEg}{dualvae_20260802-163932_ca522b}
\newcommand{\runDpp}{checkpoints/dualvae/dualvae_20260727-142050_e673a1}
\newcommand{\tagDpp}{dualvae_20260727-142050_e673a1}
\newcommand{\runVQ}{checkpoints/vqvae/vqvae_20260728-181507_9a2853}
\newcommand{\tagVQ}{vqvae_20260728-181507_9a2853}
\newcommand{\runVAE}{checkpoints/vae/vae_20260729-110353_138ede}
\newcommand{\runMCM}{checkpoints/dualvae/dualvae_20260729-174921_77fb5b}
\newcommand{\tagMCM}{dualvae_20260729-174921_77fb5b}

\newcommand{\good}[1]{\textcolor{green!55!black}{#1}}
\newcommand{\bad}[1]{\textcolor{red!70!black}{#1}}

\title{\textbf{DUALVAE Latent-Space Analysis}\\[2pt]
\large From a diagnosis of the standard prior, to a code-centered prior that\\
makes the latent a Gaussian mixture, to what that structure is worth}
\author{}
\date{}

\begin{document}
\maketitle
\vspace{-2.2em}

\begin{abstract}
\noindent
We probe a two-branch autoencoder (\textbf{DUALVAE}) that combines a discrete
vector-quantized (VQ) branch with a continuous (VAE) branch by \emph{summing} them,
$z = z_{\text{vq}} + z_{\text{cont}}$, before an attention block and the decoder. The report is
organised as one argument in five acts. \textbf{(I)}~The design promises two things: the
continuous branch does not collapse, and the sum is a \emph{Gaussian mixture} anchored on the
codewords. \textbf{(II)}~Diagnosis: with the standard $\mathcal{N}(0,I)$ prior the first claim
holds everywhere but the second fails --- the continuous spread dwarfs the codeword spacing, and
the only lever available, the KL weight $\beta$, buys mixture structure by \emph{paying}
reconstruction (a rate/separability trade-off). \textbf{(III)}~The fix: re-center the prior on the
codeword, $p(z\mid k)=\mathcal{N}(e_k,\sigma_k^2)$. This breaks the trade-off --- Runs~C/D are
statistically confirmed, well-separated mixtures at \emph{either} $\beta$, with Run~D also matching
the best reconstruction. \textbf{(IV)}~Stress tests: under a perceptual (LPIPS$+$GAN) objective all
four prior$\times\beta$ cells reconstruct \emph{identically} ($\text{rFID}\approx3.9$--$4.0$), yet
only the code-centered cells carry a mixture --- the \textbf{A+}/\textbf{B+} runs
show the GAN is \emph{not} a substitute for the prior (separability $0.24$/$1.19$, code-anchored
likelihood still negative). Pushing $\beta$ a further decade (\textbf{E+}, $\beta{=}1.0$) closes
most of the \emph{geometric} gap (separability $2.74$) but \textbf{still never flips the decisive
likelihood test} ($-0.08$), and pays for the attempt with reconstruction ($\text{rFID}\,4.94$) and
the report's only posterior collapse ($87.5\%$ active units) --- so $\beta$ and the prior are not
substitutes at any setting. A plain VQ-VAE control and a wavelet variant close off the two obvious
alternative explanations. \textbf{(V)}~Worth: the mixture makes the bag-of-codes a measurably better
linear-probe descriptor, and a matched single-branch VAE shows the dual model buys that structure at
\emph{no} reconstruction cost. Headline: \textbf{the Gaussian mixture comes from the prior's shape
--- not from $\beta$, not from the GAN, not from the codebook --- and it is free.}
\end{abstract}

\part{Claims and setup}

\section{Setup in plain words}
\label{sec:setup}
The model squeezes an image into a small grid of latent vectors (here an $8$-channel grid
of $32\times32$ locations, $256\times256$ inputs, downsample~$8$). At each location it
produces \emph{two} things and adds them:
\begin{itemize}
  \item \textbf{$z_{\text{vq}}$ (discrete):} the nearest entries from a learned codebook of
  $256$ ``words'' (with residual depth~$2$). Think of it as picking a \emph{category-like
  anchor} for that patch.
  \item \textbf{$z_{\text{cont}}$ (continuous):} a normally-distributed offset from a VAE
  branch, meant to add the leftover detail.
\end{itemize}

\textbf{The two claims under test.}
\begin{itemize}
  \item[\textbf{(A)}] The continuous branch does \emph{not} suffer posterior collapse.
  \item[\textbf{(B)}] The sum $z=z_{\text{vq}}+z_{\text{cont}}$ is a \emph{Gaussian mixture}
  --- $p(z)=\sum_k\pi_k\,\mathcal{N}(e_k,\Sigma_k)$, centers $=$ codewords --- rather than one
  smeared blob.
\end{itemize}

\textbf{The two knobs that could produce it.} The KL weight $\beta$ controls how hard the
continuous branch is pushed toward its prior; the \emph{prior shape} controls \emph{where} it is
pushed. Runs~\textbf{A} ($\beta{=}0.001$) and~\textbf{B} ($\beta{=}0.1$) use the standard
$\mathcal{N}(0,I)$ prior and isolate $\beta$. Runs~\textbf{C} ($\beta{=}0.1$) and~\textbf{D}
($\beta{=}0.001$) replace it with a \emph{code-centered} prior: the continuous branch encodes a
residual $\Delta$ around the code and the prior becomes $\mathcal{N}(e_k,\sigma_k^2)$ (flags
\texttt{residual\_continuous}, \texttt{component\_prior}, \texttt{use\_ema\_codebook}; see
\texttt{reports/code\_centered\_prior.pdf}). A ``$+$'' suffix means the same recipe retrained with
a perceptual objective (LPIPS $+$ PatchGAN), so \textbf{A+/B+/C+/D+} form a complete
$2\times2$ of \emph{prior} $\times$ $\beta$ under that objective --- the design that lets us
separate the prior's effect from the GAN's. Later runs add a wavelet detail branch (``$++$''), a
cross-view code-invariance term (CVI) and masked code modeling (MCM).

All numbers below come from $200$ validation images (latent statistics), the \textbf{full
validation set ($3925$ images)} (reconstruction, FID/KID, probes), $60$ images (code
segmentation) and $5$ images (feature pictures). Table~\ref{tab:ckpts} maps every run label to
its checkpoint directory.

\begin{table}[h]
\centering
\caption{Run label $\to$ checkpoint directory (all under
\texttt{checkpoints/dualvae/} unless noted). ``code'' = code-centered prior; ``+'' = LPIPS+GAN;
``++'' = + wavelet; CVI = + cross-view code-invariance; MCM = + masked code modeling.
The four \textbf{bold} rows are the $2\times2$ (prior $\times$ $\beta$) under the perceptual
objective.}
\label{tab:ckpts}
\small
\begin{tabular}{lll}
\toprule
Run & Recipe & Checkpoint directory\\
\midrule
A     & $\beta{=}.001$, standard prior      & \texttt{dualvae\_20260724-121508\_392b5b}\\
B     & $\beta{=}.1$, standard prior        & \texttt{dualvae\_20260724-133147\_b03f09}\\
C     & $\beta{=}.1$, code prior            & \texttt{dualvae\_20260725-015430\_5878cc}\\
D     & $\beta{=}.001$, code prior          & \texttt{dualvae\_20260725-134828\_ce6208}\\
\midrule
\textbf{A+} & A + LPIPS/GAN (std prior, $\beta{=}.001$) & \texttt{dualvae\_20260801-201243\_0b2ee6}\\
\textbf{B+} & B + LPIPS/GAN (std prior, $\beta{=}.1$)   & \texttt{dualvae\_20260801-201309\_1e27ab}\\
\textbf{C+} & C + LPIPS/GAN (code prior, $\beta{=}.1$)  & \texttt{dualvae\_20260726-230719\_5e588d}\\
\textbf{D+} & D + LPIPS/GAN (code prior, $\beta{=}.001$)& \texttt{dualvae\_20260726-230704\_866e33}\\
\midrule
E+    & LPIPS/GAN, std prior, $\beta{=}1.0$ (the $\beta$ ceiling) & \texttt{dualvae\_20260802-163932\_ca522b}\\
\midrule
C++   & C+ + wavelet                        & \texttt{dualvae\_20260727-142016\_2f1dfc}\\
D++   & D+ + wavelet                        & \texttt{dualvae\_20260727-142050\_e673a1}\\
D+CVI & D+ + code-invariance                & \texttt{dualvae\_20260727-210658\_f2145f}\\
D+MCM & D+ + masked code modeling           & \texttt{dualvae\_20260729-174921\_77fb5b}\\
VQ-VAE & codes only, $+$LPIPS/GAN (no cont.\ branch) & \texttt{vqvae/vqvae\_20260728-181507\_9a2853}\\
VAE   & continuous only, matched (8-ch, $+$LPIPS/GAN) & \texttt{vae/vae\_20260729-110353\_138ede}\\
\bottomrule
\end{tabular}
\end{table}

\part{Diagnosis --- what the standard prior actually learns}

\section{Claim A holds --- everywhere except at the $\beta$ ceiling}
\label{sec:collapse}
\textbf{What ``collapse'' means.} A VAE branch collapses when its output stops depending on
the image: the posterior $q(z\mid x)$ just becomes the prior $\mathcal{N}(0,I)$ (variance
$\to 1$, mean $\to 0$), so the latent carries no information and the decoder ignores it.

\textbf{What we measured (and why it is easy to read).}
\begin{itemize}
  \item \textbf{Active units:} count latent dimensions that actually change from image to
  image ($\mathrm{Var}_x(\mathbb{E}[z\mid x])>0.01$). Collapse $\Rightarrow$ near $0$.
  \item \textbf{KL to the prior (nats):} how far the posterior is from the prior.
  Collapse $\Rightarrow$ near $0$.
  \item \textbf{Posterior std / $|$mean$|$:} collapse sits at std$=1$, mean$=0$.
\end{itemize}

\begin{table}[h]
\centering
\caption{Collapse diagnostics (continuous branch). Every run in the original design --- both
priors, both $\beta$, with and without the GAN --- is far from the collapse values, so
\textbf{none of them collapses}. Runs~A/D and the four $2\times2$ GAN runs are near-deterministic
and high-rate; Runs~B/C are more regularized but still healthy. The exception is
\textbf{E+} ($\beta{=}1.0$), added later: one channel of eight is switched off entirely and the
posterior std climbs back toward the prior --- see \S\ref{sec:betaceiling}.
``code'' = code-centered prior, ``+'' = LPIPS/GAN.}
\label{tab:collapse}
\small
\setlength{\tabcolsep}{4.5pt}
\begin{tabular}{lcccccccc}
\toprule
Metric & Collapse & \textbf{Run~A} & \textbf{Run~B} & \textbf{Run~C} & \textbf{Run~D} & \textbf{A+} & \textbf{B+} & \textbf{E+}\\
       & looks like & ($\beta{=}.001$) & ($\beta{=}.1$) & (code,$.1$) & (code,$.001$) & (std,$.001$) & (std,$.1$) & (std,$1.0$)\\
\midrule
Active units            & $\approx 0$ & \good{$8192/8192$} & \good{$6144/8192$} & \good{$8192/8192$} & \good{$8192/8192$} & \good{$8192/8192$} & \good{$8192/8192$} & \bad{$7168/8192$}\\
KL nats / element       & $\approx 0$ & \good{$4.033$}     & \good{$0.594$}     & \good{$1.042$}     & \good{$2.574$}     & \good{$7.456$}    & \good{$2.444$} & \bad{$0.323$}\\
Posterior std (avg)     & $1.0$       & \good{$0.022$}     & \good{$0.616$}     & \good{$0.224$}     & \good{$0.047$}     & \good{$0.006$}    & \good{$0.103$} & \bad{$0.673$}\\
$|$Posterior mean$|$ avg& $0.0$       & \good{$0.800$}     & \good{$0.356$}     & \good{$0.140$}     & \good{$0.134$}     & \good{$1.956$}    & \good{$0.809$} & $0.391$\\
\bottomrule
\end{tabular}
\end{table}

\textbf{Plain-words result.} Every latent dimension is used and the posterior sits nowhere
near the prior, so \textbf{Claim~A holds in every run we ever trained}. A nuance worth
noting: in Run~A the posterior std is tiny ($0.022$), i.e.\ the branch became almost
\emph{deterministic} and packs in a lot of information --- which is exactly why Run~A
reconstructs so well (\S\ref{sec:branches}). The perceptual objective pushes this further still
(A+: std $0.006$, KL $7.5$ nats). Supporting plots:
\texttt{collapse\_kl\_per\_channel}, \texttt{collapse\_active\_units},
\texttt{collapse\_posterior\_stats} in each run's folder.

\textbf{The one exception, and what it costs.} Claim~A holds across the entire original design,
but it is not unconditional. \textbf{E+} ($\beta{=}1.0$, \S\ref{sec:betaceiling}) --- the only run
trained with KL pressure an order of magnitude above B+ --- drops to $7168/8192$ active units
because a single channel dies outright (per-channel KL $0.0009$ nats, active fraction $0$), and its
posterior std rebounds to $0.673$, ten times the other GAN runs. So the honest statement is:
\emph{the continuous branch does not collapse at any $\beta$ anyone would want to use, and starts
to at $\beta=1.0$ --- which is also the point where reconstruction breaks.} With that boundary
marked, Claim~A is settled; the rest of the report is about Claim~B.

\section{Claim B fails: the sum is one blob, and $\beta$ only trades it against fidelity}
\label{sec:mixture}
\textbf{The idea in one line.} Because $z = e_k + \Delta$ (a codeword plus a Gaussian
offset), the overall distribution is, by construction,
$p(z)=\sum_k \pi_k\,\mathcal{N}(e_k,\Sigma_k)$: a Gaussian mixture whose \emph{centers are
the codewords} and whose \emph{weights are how often each code is used}. The real question
is whether those components are \emph{separated} or just smeared into one blob.

\textbf{What we measured.}
\begin{itemize}
  \item \textbf{Separability ratio} $=$ (typical distance between code centers) $/$
  (continuous spread). $>1$ means the codes stand apart from the noise; $<1$ means the
  continuous offset drowns them.
  \item \textbf{Single-Gaussian vs.\ mixture:} BIC (model selection) and the held-out
  log-likelihood of the \emph{code-anchored} mixture vs.\ a single Gaussian. The latter is the
  decisive test: it fixes the component means \emph{to the codewords} and asks whether that
  explains the latent better than one Gaussian --- exactly Claim~B.
  \item \textbf{Cluster quality:} silhouette using the code as the label.
  \item \textbf{Multimodality:} project onto the axis between the two \emph{most-separated}
  code centers; report a 2-mode separation (distance between the two fitted modes in units
  of their width) and Hartigan's dip test.
\end{itemize}

\begin{table}[h]
\centering
\caption{Gaussian-mixture diagnostics for $z=z_{\text{vq}}+z_{\text{cont}}$ under the
\textbf{standard} prior. Run~A is one smeared blob; Run~B is \emph{emerging} but not confirmed.
The perceptual objective (A+/B+, analysed in full in \S\ref{sec:twobytwo}) does \emph{not}
rescue it. Only the code-centered prior (Runs~C/D, \S\ref{sec:codeprior}) passes every test
--- shown here for contrast.}
\label{tab:gmm}
\small
\setlength{\tabcolsep}{4pt}
\begin{tabular}{lcccccc}
\toprule
 & \multicolumn{4}{c}{standard prior} & \multicolumn{2}{c}{code-centered prior}\\
\cmidrule(r){2-5}\cmidrule(l){6-7}
Metric (want for a clean GMM) & \textbf{A} & \textbf{B} & \textbf{A+} & \textbf{B+} & \textbf{C} & \textbf{D}\\
 & ($.001$) & ($.1$) & ($.001$,GAN) & ($.1$,GAN) & ($.1$) & ($.001$)\\
\midrule
Codes used                       & $194/256$ & $255/256$ & \bad{$101/256$} & $190/256$ & \good{$256/256$} & \good{$256/256$}\\
Codeword spacing                 & $0.108$   & $0.643$   & $0.576$   & $1.325$   & $1.121$ & $0.583$\\
Continuous (residual) std        & $1.005$   & $0.948$   & $2.398$   & $1.110$   & $0.412$ & $0.248$\\
\textbf{Separability ratio} ($\gg 1$) & \bad{$0.11$} & $0.68$ & \bad{$0.24$} & $1.19$ & \good{$2.72$} & \good{$2.35$}\\
2-mode separation ($\gg 1$)      & \bad{$0.31$} & \good{$2.09$} & $1.51$ & $1.70$ & \good{$1.96$} & \good{$1.94$}\\
\textbf{Code-anchored vs.\ single LL} ($>0$) & \bad{$-3.11$} & $-0.09$ & \bad{$-3.12$} & \bad{$-0.18$} & \good{$+1.27$} & \good{$+0.63$}\\
Silhouette by code ($>0$)        & \bad{$-0.12$} & $0.006$ & \bad{$-0.06$} & $0.039$ & \good{$+0.13$} & \good{$+0.16$}\\
GMM vs.\ single-Gaussian BIC ($<0$) & \bad{$+1987$} & \bad{$+25060$} & \bad{$+14839$} & \bad{$+23414$} & \good{$-10347$} & \good{$-3508$}\\
Hartigan dip $p$ ($<0.05$)       & \bad{$0.999$} & \bad{$0.991$} & \bad{$0.999$} & \bad{$1.000$} & \good{$0.0007$} & \good{$0.005$}\\
Effective \# components          & $148$ & $146$ & $87$ & $147$ & $187$ & $194$\\
\bottomrule
\end{tabular}
\end{table}

\textbf{Plain-words result.} In Run~A the continuous spread ($\approx 1.0$) is about
$10\times$ the codeword spacing ($\approx 0.1$): the codes barely move the latent, so the
mixture is invisible --- one Gaussian blob (Fig.~\ref{fig:gmm}, left). Stronger KL pressure
in Run~B pushes the codes $6\times$ farther apart, and the picture changes: the code centers
now spread across the plane with points organizing around them (Fig.~\ref{fig:gmm}, right),
the two most-separated codes sit $2.09$ widths apart, the code-anchored mixture now ties a
single Gaussian, and the silhouette climbs out of negative. It is a mixture
\emph{coming into focus}, though the strict whole-distribution tests (BIC, dip) have not
flipped yet, because $\sim\!190$ overlapping components still fill a smooth 1-D marginal.

\textbf{And the improvement is not free.} That is the whole problem with $\beta$ as the knob:
Run~B buys its emerging mixture with reconstruction (PSNR $29.4\to27.3$, rFID $24.4\to40.0$;
Table~\ref{tab:ablation}). Under the standard prior, mixture structure and fidelity are in
\emph{competition} --- a \textbf{rate/separability trade-off}. Everything in
Part~III is about escaping it.

\begin{figure}[h]
\centering
\begin{subfigure}{0.49\textwidth}
  \includegraphics[width=\linewidth]{\runA/gmm_pca_codes_\tagA.pdf}
  \caption{Run~A ($\beta{=}0.001$): centers piled at the origin, colors mixed.}
\end{subfigure}\hfill
\begin{subfigure}{0.49\textwidth}
  \includegraphics[width=\linewidth]{\runB/gmm_pca_codes_\tagB.pdf}
  \caption{Run~B ($\beta{=}0.1$): centers spread out, points organize by code.}
\end{subfigure}
\caption{Per-location latent $z=z_{\text{vq}}+z_{\text{cont}}$ (PCA to 2-D), colored by code,
$\bigstar =$ code centers. Stronger $\beta$ turns one blob into an emerging mixture --- at a
reconstruction cost.}
\label{fig:gmm}
\end{figure}

\section{Who carries the image: feature maps and branch ablation}
\label{sec:branches}
If the codes are invisible in the latent, what are they \emph{doing}? Two views answer this: a
picture of each branch's features, and a measurement of what happens when a branch is removed.

\subsection*{Feature maps (the DINO trick)}
For a few images we take each branch's latent grid, run PCA over the location vectors, and paint
the top-3 components as an RGB image (and the 1st component as a heatmap). Similar features get
similar colors, so we can literally \emph{see} what each branch encodes, painted back onto the
image.

\textbf{Plain-words result} (Figs.~\ref{fig:featA} and~\ref{fig:featB}). Reconstructions are
nearly identical to the originals in both runs. The \textbf{VQ features} trace object shapes
and layout (the truck, the headphones, the gas pump are visible as coherent colored
regions); the \textbf{continuous features} look more like fine, high-frequency texture; the
\textbf{summed PC1 heatmap} lights up object contours, just like DINO's ``foreground''
component. So the discrete branch carries \emph{structure} and the continuous branch carries
\emph{detail} --- consistent with the ablation below. The effect is visible in both runs;
in Run~B (higher $\beta$) the VQ features carry even more of the object structure, matching
the ablation finding that the codes do more of the reconstruction work there. (Equivalent
figures for every other run, including A+/B+, are in each run's folder.)

\begin{figure}[h]
\centering
\includegraphics[width=0.92\textwidth]{\runA/feature_pca_images_\tagA.pdf}
\caption{\textbf{Run~A} ($\beta{=}0.001$) latent features projected onto $5$ validation
images (PCA$\to$RGB). Columns: original, reconstruction, VQ features, continuous features,
summed features, summed PC1.}
\label{fig:featA}
\end{figure}

\begin{figure}[h]
\centering
\includegraphics[width=0.92\textwidth]{\runB/feature_pca_images_\tagB.pdf}
\caption{\textbf{Run~B} ($\beta{=}0.1$) latent features projected onto the same $5$
validation images (PCA$\to$RGB), same column order. Compare the VQ-features column against
Run~A (Fig.~\ref{fig:featA}).}
\label{fig:featB}
\end{figure}

\subsection*{Branch ablation}
We reconstruct the \textbf{full validation set ($3925$ images)} three
ways: \emph{full} (both branches), \emph{vq\_only} (continuous zeroed), and \emph{vanilla\_only}
(VQ zeroed), and measure PSNR/SSIM (higher is better) and LPIPS (lower is better). If a branch
had collapsed, deleting it would change nothing; a big drop means the branch is load-bearing.
\emph{(All reconstruction/FID numbers in this report are computed on the full val set; the
earlier $1500$-image FIDs were biased upward $\sim$$2\times$ by the small sample.)}

\begin{table}[h]
\centering
\caption{Reconstruction quality by branch ablation (\textbf{full val, $3925$
images}). The roles flip between Run~A and Run~B; in \textbf{Runs~C/D the codes carry almost
everything} (\texttt{vq\_only}$\approx$\texttt{full}). \textbf{Run~D} matches Run~A's fidelity,
has a clean mixture, \emph{and} has the best \textbf{rFID} of the MSE runs. The two GAN rows
(A+/B+) are included for completeness and discussed in \S\ref{sec:twobytwo}; rFID is shown only
for \texttt{full}.}
\label{tab:ablation}
\begin{tabular}{llcccc}
\toprule
Run & Mode & PSNR\,$\uparrow$ & SSIM\,$\uparrow$ & LPIPS\,$\downarrow$ & rFID\,$\downarrow$\\
\midrule
\multirow{3}{*}{\textbf{A} ($\beta{=}0.001$)}
 & full          & $29.37$ & $0.878$ & $0.112$ & $24.4$\\
 & vq\_only      & \bad{$10.33$} & \bad{$0.323$} & \bad{$0.850$} & --\\
 & vanilla\_only & $21.56$ & $0.787$ & $0.203$ & --\\
\midrule
\multirow{3}{*}{\textbf{B} ($\beta{=}0.1$)}
 & full          & $27.29$ & $0.814$ & $0.189$ & \bad{$40.0$}\\
 & vq\_only      & $19.93$ & $0.614$ & $0.465$ & --\\
 & vanilla\_only & \bad{$10.68$} & \bad{$0.326$} & \bad{$0.831$} & --\\
\midrule
\multirow{3}{*}{\textbf{C} (code, $\beta{=}0.1$)}
 & full          & $27.98$ & $0.835$ & $0.172$ & $31.9$\\
 & vq\_only      & \good{$26.92$} & \good{$0.803$} & \good{$0.219$} & --\\
 & vanilla\_only & \bad{$10.73$} & \bad{$0.278$} & \bad{$0.835$} & --\\
\midrule
\multirow{3}{*}{\textbf{D} (code, $\beta{=}0.001$)}
 & full          & \good{$28.89$} & \good{$0.869$} & \good{$0.117$} & \good{$22.7$}\\
 & vq\_only      & \good{$26.45$} & \good{$0.806$} & \good{$0.199$} & --\\
 & vanilla\_only & \bad{$10.78$} & \bad{$0.384$} & \bad{$0.780$} & --\\
\midrule
\multirow{3}{*}{\textbf{A+} (std, $\beta{=}0.001$, GAN)}
 & full          & $27.26$ & $0.834$ & \good{$0.043$} & \good{$4.02$}\\
 & vq\_only      & \bad{$8.24$} & \bad{$0.328$} & \bad{$0.780$} & --\\
 & vanilla\_only & $14.09$ & $0.527$ & $0.245$ & --\\
\midrule
\multirow{3}{*}{\textbf{B+} (std, $\beta{=}0.1$, GAN)}
 & full          & $26.94$ & $0.829$ & \good{$0.043$} & \good{$3.70$}\\
 & vq\_only      & $18.85$ & $0.548$ & $0.586$ & --\\
 & vanilla\_only & \bad{$10.91$} & \bad{$0.449$} & $0.410$ & --\\
\bottomrule
\end{tabular}
\end{table}

\textbf{Plain-words result.} In \textbf{Run~A} the \emph{continuous} branch does almost all
the work: alone it reaches PSNR $21.6$, while codes-alone is near-useless ($10.3$) --- the
continuous branch adds $+19$ PSNR. In \textbf{Run~B} the roles \emph{reverse}: the
\emph{codes} now carry the image (PSNR $19.9$ alone) and the continuous branch adds detail
on top ($+7.3$ PSNR). Deleting either branch hurts a lot in both runs, so \emph{both
branches are used} --- another confirmation of Claim~A. On \textbf{rFID} (the perceptual
metric, full model, full val) the ordering among the MSE runs is $\textbf{D}\,22.7 <
\textbf{A}\,24.4 < \textbf{C}\,31.9 < \textbf{B}\,40.0$: the code-prior Run~D is the most
realistic --- slightly ahead of Run~A despite Run~A's higher PSNR --- while $\beta{=}0.1$ without
the code prior (Run~B) is the least realistic. See
\texttt{ablation\_grid\_*} and \texttt{ablation\_metrics\_*}; Fig.~\ref{fig:abl} shows the
Run~A grid, where \texttt{vq\_only} is visibly noise.

\begin{figure}[h]
\centering
\includegraphics[width=0.62\textwidth]{\runA/ablation_grid_\tagA.pdf}
\caption{Run~A branch-ablation reconstructions. Codes-only (\texttt{vq\_only}) is noise;
continuous-only (\texttt{vanilla\_only}) already looks almost like the full model --- the
continuous branch is doing the heavy lifting here.}
\label{fig:abl}
\end{figure}

\section{What the codes \emph{really} contribute --- beyond zeroing}
\label{sec:codeinfo}
\S\ref{sec:branches} shows that deleting a branch hurts, but it has a hidden confound: the decoder
was trained on the \emph{sum} $z_{\text{vq}}+z_{\text{cont}}$, so zeroing a branch removes
information \emph{and} pushes the decoder input off the manifold it ever saw. A direct
measurement makes this concrete: in Run~A most of the discrete signal is a single
\textbf{constant offset} added at every location (shared-DC part $\|\overline{z_{\text{vq}}}\|
= 0.79$, which is $64\%$ of $\|z_{\text{vq}}\|$); in Run~B the codes are mostly
location-specific ($82\%$). So ``zeroing the codes'' partly just deletes a constant the
decoder leans on. We ran four experiments to pull these effects apart.

\paragraph{Replace-with-mean and random-code swap (in easy words).}
Keeping the continuous branch fixed, we feed the decoder four versions of the discrete part:
the \emph{real} codes (\texttt{full}); the \emph{dataset-average} code at every location
(\texttt{code\_mean} --- keeps the constant offset, throws away per-image code information);
the real codes \emph{shuffled} across locations (\texttt{code\_random} --- same code
statistics, wrong image); and \emph{zero} (\texttt{code\_zero}, i.e.\ \texttt{vanilla\_only}).
This lets us split the ``codes help'' gap into two named pieces:
\[
\underbrace{Q(\texttt{full})-Q(\texttt{code\_mean})}_{\text{real per-image code INFO}}
\quad\text{and}\quad
\underbrace{Q(\texttt{code\_mean})-Q(\texttt{code\_zero})}_{\text{DC / off-manifold effect}} .
\]

\begin{table}[h]
\centering
\caption{Code-intervention reconstruction (SSIM$\uparrow$ / PSNR$\uparrow$ / LPIPS$\downarrow$)
and the decomposition of the \texttt{full}$-$\texttt{vanilla\_only} gap. \textbf{The reason
the codes ``help'' differs sharply: a DC artifact in Run~A, but genuine per-image information
in Runs~B/C/D.}}
\label{tab:codeinfo}
\scriptsize
\setlength{\tabcolsep}{3pt}
\begin{tabular}{lcccccc}
\toprule
 & \textbf{A} ($\beta{=}.001$) & \textbf{B} ($\beta{=}.1$) & \textbf{C} (code,$.1$) & \textbf{D} (code,$.001$) & \textbf{A+} (std,GAN) & \textbf{B+} (std,GAN)\\
\midrule
\texttt{full}         & $0.882 / 29.27 / 0.109$ & $0.819 / 27.19 / 0.185$ & $0.839 / 27.88 / 0.168$ & $0.873 / 28.77 / 0.113$ & $0.839 / 27.05 / 0.043$ & $0.834 / 26.72 / 0.043$\\
\texttt{code\_mean}   & $0.849 / 27.52 / 0.135$ & $0.339 / 11.54 / 0.805$ & $0.293 / 10.62 / 0.836$ & $0.355 / 10.76 / 0.772$ & $0.780 / 24.11 / 0.088$ & $0.468 / 11.74 / 0.420$\\
\texttt{code\_random} & $0.825 / 26.71 / 0.153$ & $0.231 / 11.56 / 0.606$ & $0.123 / 9.11 / 0.760$ & $0.151 / 9.55 / 0.719$ & $0.759 / 23.53 / 0.107$ & $0.397 / 11.12 / 0.427$\\
\texttt{code\_zero}   & $0.791 / 21.59 / 0.198$ & $0.331 / 10.62 / 0.830$ & $0.282 / 10.56 / 0.837$ & $0.386 / 10.47 / 0.782$ & $0.539 / 14.08 / 0.240$ & $0.456 / 10.85 / 0.409$\\
\midrule
real code INFO (\texttt{full}$-$\texttt{code\_mean}) & \good{$+0.033$} & \good{$+0.480$} & \good{$+0.546$} & \good{$+0.518$} & \good{$+0.059$} & \good{$+0.366$}\\
DC effect (\texttt{code\_mean}$-$\texttt{code\_zero}) & \bad{$+0.058$} & $+0.008$ & $+0.011$ & $-0.031$ & \bad{$+0.241$} & $+0.013$\\
share of gap that is real info & $\approx 36\%$ & $\approx 98\%$ & $\approx 98\%$ & $\approx 100\%$ & \bad{$\approx 20\%$} & $\approx 97\%$\\
\bottomrule
\end{tabular}
\end{table}

\textbf{Plain-words result.} In \textbf{Run~A} about two-thirds of the ``codes help'' gap is
just the decoder missing a constant offset (an \emph{artifact} of zeroing, not information);
only $\sim$one-third is real per-image code content --- visible in the grid, where
\texttt{code\_mean} looks almost like \texttt{full} while \texttt{code\_zero} is
\emph{desaturated}. In \textbf{Run~B} it flips: $\sim$$98\%$ of the gap is genuine per-image
information (dropping the codes' spatial pattern destroys the image), and the constant offset
barely matters. The random-code swap agrees: shuffling the codes costs $+0.06$ SSIM in Run~A
but a huge $+0.59$ in Run~B, i.e.\ putting the \emph{right} code in the \emph{right} place
matters enormously in Run~B and only a little in Run~A. \textbf{This is the second face of the
same trade-off:} at low $\beta$ under the standard prior the codes are nearly decorative.
The perceptual objective does not repair this either --- it makes it \emph{worse}: in \textbf{A+}
only $\approx$$20\%$ of the gap is real per-image information (vs Run~A's $36\%$), the lowest of
any run, matching the half-dead codebook of \S\ref{sec:twobytwo}; \textbf{B+} sits at
$\approx$$97\%$, like B. The share of genuine code information tracks the same thing everything
else does --- whether the continuous branch is constrained relative to the codes.

\paragraph{Which frequencies each branch carries (Fig.~\ref{fig:spec}).}
We take the pixel-space contribution of each branch ($\Delta$codes $=$
\texttt{full}$-$\texttt{code\_mean}, $\Delta$cont $=$ \texttt{full}$-$\texttt{cont\_zero})
and look at its power across spatial frequencies. In Run~A the continuous branch dominates
\emph{every} frequency (it does all the work). In Run~B the curves cross: the \textbf{codes
carry the low/coarse frequencies} (the bulk of the image) while the \textbf{continuous branch
takes over at high/fine frequencies} (the texture) --- a clean coarse/detail division of
labour.

\begin{figure}[h]
\centering
\begin{subfigure}{0.49\textwidth}
  \includegraphics[width=\linewidth]{\runA/branch_contribution_spectrum_\tagA.pdf}
  \caption{Run~A: continuous (blue) dominates all bands.}
\end{subfigure}\hfill
\begin{subfigure}{0.49\textwidth}
  \includegraphics[width=\linewidth]{\runB/branch_contribution_spectrum_\tagB.pdf}
  \caption{Run~B: codes (red) carry low freq, continuous takes high freq.}
\end{subfigure}
\caption{Radially-averaged power of each branch's pixel-space contribution.}
\label{fig:spec}
\end{figure}

\paragraph{A fair standalone-VAE baseline.} The correct test of ``do the codes
improve a VAE'' is not \texttt{vanilla\_only} (that branch was \emph{trained} to lean on the
codes) but a matched continuous-only VAE trained from scratch. We use a fully matched one
(\texttt{vae/vae\_20260729-110353\_138ede}): same $8$-channel latent, same LPIPS+GAN recipe,
$\beta{=}0.001$, $200$ epochs --- i.e.\ \emph{exactly Run~D+ minus the VQ branch and code-centered
prior}. The key result: \textbf{each dual model's \texttt{vanilla\_only} ($0.791$, $0.331$, $0.279$,
$0.382$ SSIM) is far \emph{worse} than this standalone VAE} ($0.832$), confirming that the dual's
continuous branch offloaded coarse structure to the codes during training --- which is exactly why
the plain branch ablation over-credits the codes. What this baseline settles about
\emph{reconstruction} is taken up in \S\ref{sec:gan}. \emph{(This supersedes the earlier
$4$-channel MSE VAE baseline and its capacity caveat.)}

\part{The fix --- a code-centered prior}

\section{Moving the prior onto the code (Runs~C and D)}
\label{sec:codeprior}
\emph{Run~C uses the code-centered prior with $\beta=0.1$; Run~D uses the identical recipe with
$\beta=0.001$. Together they show the mixture comes from the prior \emph{shape}, not the KL
strength --- and that low $\beta$ (Run~D) buys back the reconstruction quality.}

\subsection*{Why the standard prior cannot make a mixture}
In Runs~A and B the continuous branch $z_{\text{cont}}=\mu+\sigma\odot\varepsilon$ is
regularized by $\mathrm{KL}\!\big(\mathcal{N}(\mu,\sigma^2)\,\|\,\mathcal{N}(0,I)\big)$, which
pushes $\mu\to 0$ --- toward the \emph{origin}, and \emph{independently of the code}. The code
$z_{\text{vq}}=e_k$ is merely added afterwards. So the continuous branch is free to fill a
full unit-variance Gaussian whose spread ($\sim\!1.0$) dwarfs the codeword spacing
($\sim\!0.1$): the codewords are tiny bumps inside one big cloud and no mixture is visible.
This is not fixed by ``summing the code into the mean before the reparameterization'':
$z_{\text{vq}}+(\mu+\sigma\varepsilon)$ and $(\mu+z_{\text{vq}})+\sigma\varepsilon$ are the
\emph{same} number, so with the same KL nothing changes.

\subsection*{The fix: make the codeword the \emph{center} of the prior}
Write the continuous branch as a \emph{residual} $\Delta$ around the code, $z=e_k+\Delta$ with
$\Delta=\mu_\Delta+\sigma\odot\varepsilon$, and change the prior so that, \emph{given} the
code $k$, the latent is expected to sit \emph{on} the codeword:
\begin{equation}
p(z\mid k)=\mathcal{N}(e_k,\sigma_k^2 I)
\quad\Longleftrightarrow\quad
p(\Delta\mid k)=\mathcal{N}(0,\sigma_k^2 I).
\end{equation}
The regularizer becomes a KL between the residual posterior and this code-centered prior,
which has a clean per-dimension form:
\begin{equation}
\mathrm{KL}\!\big(\mathcal{N}(\mu_\Delta,\sigma^2)\,\|\,\mathcal{N}(0,\sigma_k^2)\big)
=\log\frac{\sigma_k}{\sigma}+\frac{\sigma^2+\mu_\Delta^2}{2\sigma_k^2}-\tfrac12 .
\label{eq:klc}
\end{equation}
The decisive term is $\mu_\Delta^2/(2\sigma_k^2)$: it penalizes the residual \emph{mean}, so
training now drives $\mu_\Delta\to 0$, i.e.\ it pulls $z$ \emph{toward the codeword $e_k$}
instead of toward the origin. Averaging over how often each code is used ($\pi_k$) gives the
marginal prior
\begin{equation}
\boxed{\,p(z)=\sum_k \pi_k\,\mathcal{N}(e_k,\sigma_k^2 I)\,}
\end{equation}
--- a genuine \textbf{Gaussian mixture} whose component \emph{means are the codewords}, weights
are the code frequencies, and widths are the per-code $\sigma_k$ (estimated online as the EMA
within-code variance $v_k/N_k$). The components are well separated exactly when the within-code
spread is smaller than the gap between codewords,
\begin{equation}
\text{separated}\ \Longleftrightarrow\ \sigma_k\ \ll\ \min_{j\neq k}\lVert e_j-e_k\rVert ,
\label{eq:sepc}
\end{equation}
which \eqref{eq:klc} now actively enforces (small residual) while the codes drift apart to carry
information. In the code this is \texttt{residual\_continuous} + \texttt{component\_prior} +
\texttt{use\_ema\_codebook} (full derivation in \texttt{reports/code\_centered\_prior.pdf}).

\textbf{It works --- decisively.} Every Claim-B metric flips to the mixture side
(Table~\ref{tab:gmm}): the continuous residual std \emph{drops} from $\approx\!1.0$ to $0.41$
while the codeword spacing \emph{grows} to $1.12$, so the \textbf{separability ratio jumps to
$2.72$} (codes now stand well apart from the within-code noise, the condition $\sigma_k\ll$
spacing). For the first time the K-component GMM \textbf{beats} a single Gaussian on BIC
($-10347$), the code-anchored mixture \textbf{beats} a single Gaussian in held-out likelihood
($+1.27$), the silhouette is positive ($+0.13$), and \textbf{Hartigan's dip test is significant}
($p=7\times10^{-4}$, multimodal). The PCA view (Fig.~\ref{fig:gmmC}) shows what the numbers mean:
tight, separated clusters, each sitting on its codeword --- compare the single blob of Run~A.

\textbf{And it is nearly free.} Run~C keeps the continuous branch healthy (no collapse:
$100\%$ active units, KL $1.04$), and its reconstruction ($0.836$ SSIM / $27.85$ PSNR) is
\emph{better than Run~B} ($0.819$ / $27.11$) and close to Run~A --- so we gained a clean,
separated mixture \emph{without} paying the reconstruction cost the plain $\beta$ knob demanded.
The ablation (Table~\ref{tab:ablation}) confirms the intended division of labour: the codes now
carry almost the entire image (\texttt{vq\_only} $0.805\approx$ \texttt{full} $0.836$) and the
continuous branch is a small refinement ($+0.03$ SSIM), and the decomposition shows
$\approx\!98\%$ of the code contribution is genuine per-image information
(Table~\ref{tab:codeinfo}).

\begin{figure}[h]
\centering
\includegraphics[width=0.62\textwidth]{\runC/gmm_pca_codes_\tagC.pdf}
\caption{Run~C per-location latent $z=z_{\text{vq}}+z_{\text{cont}}$ (PCA to 2-D), colored by
code, $\bigstar =$ code centers. The code-centered prior produces tight, well-separated clusters
sitting on their codewords --- a genuine Gaussian mixture (separability $2.72$, dip $p<10^{-3}$).
Contrast Run~A (Fig.~\ref{fig:gmm}a, one blob).}
\label{fig:gmmC}
\end{figure}

\textbf{Feature maps and frequency (Run~C).} The DINO-style feature maps (Fig.~\ref{fig:featC})
show the same reconstruction quality with even cleaner object structure in the VQ column, and
the frequency decomposition (Fig.~\ref{fig:specC}) makes the new division of labour explicit:
the codes ($\Delta$codes) now dominate at \emph{every} frequency (they carry the image, about
$1000\times$ the continuous power) while the continuous branch ($\Delta$cont) adds only a thin
correction, relatively stronger toward high/fine frequencies --- the signature of ``codes =
mixture-component identity, continuous = small within-component detail''.

\begin{figure}[h]
\centering
\includegraphics[width=0.92\textwidth]{\runC/feature_pca_images_\tagC.pdf}
\caption{\textbf{Run~C} latent features projected onto $5$ validation images (PCA$\to$RGB),
same column order as Figs.~\ref{fig:featA}--\ref{fig:featB}: original, reconstruction, VQ
features, continuous features, summed features, summed PC1.}
\label{fig:featC}
\end{figure}

\begin{figure}[h]
\centering
\includegraphics[width=0.62\textwidth]{\runC/branch_contribution_spectrum_\tagC.pdf}
\caption{\textbf{Run~C} radially-averaged power of each branch's pixel-space contribution.
The codes (red) dominate at all frequencies (they carry the image, $\sim\!1000\times$ the
continuous power); the continuous branch (blue) adds only a small correction, relatively
stronger at high/fine frequencies --- the clean ``codes carry the image, continuous refines it''
split the mixture wiring is designed to produce.}
\label{fig:specC}
\end{figure}

\subsection*{Run~D: the mixture is free at low $\beta$ too --- and reconstruction returns}
Run~D keeps the code-centered prior but drops $\beta$ from $0.1$ to $0.001$. The mixture
\textbf{survives essentially unchanged}: separability $2.35$, GMM still beats a single Gaussian
on BIC ($-3508$) and held-out likelihood ($+0.63$), silhouette $+0.16$, dip test still
significant ($p=0.005$, multimodal), all $256$ codes used (Fig.~\ref{fig:gmmD}). So the mixture
is produced by the prior's \emph{shape} (the codeword-centered mean), \emph{not} by the strength
of the KL --- exactly what \eqref{eq:klc} predicts, since the KL only shrinks the residual, and
even a weak KL keeps it below the codeword spacing ($0.25 < 0.58$).

\textbf{What the low $\beta$ buys is reconstruction.} With the mixture no longer depending on a
strong KL, the weaker $\beta$ lets the continuous branch run at higher rate (KL $2.57$ nats,
posterior std $0.047$ --- near-deterministic, like Run~A), and reconstruction climbs to
\textbf{$0.870$ SSIM / $28.73$ PSNR}, essentially matching Run~A ($0.882 / 29.09$) and clearly
above Run~C ($0.836 / 27.85$). Run~D therefore \textbf{breaks the rate/separability trade-off
outright}: it is simultaneously (i) a clean, confirmed Gaussian mixture and (ii) an
almost-best-in-class reconstructor, with no collapse ($100\%$ active units) and $\approx\!100\%$
of the code contribution being genuine per-image information (Table~\ref{tab:codeinfo}). The only
price relative to Run~C is slightly tighter, less widely-spaced components (spacing $0.58$ vs
$1.12$), which cost nothing on any mixture test.

\begin{figure}[h]
\centering
\includegraphics[width=0.62\textwidth]{\runD/gmm_pca_codes_\tagD.pdf}
\caption{\textbf{Run~D} (code-centered prior, $\beta{=}0.001$) per-location latent, PCA to 2-D,
colored by code, $\bigstar =$ code centers. Same well-separated mixture as Run~C
(Fig.~\ref{fig:gmmC}) --- just more compact --- but with Run~A-level reconstruction.}
\label{fig:gmmD}
\end{figure}

\part{Stress tests --- does the mixture survive a real objective?}

\section{Adding a perceptual (LPIPS + GAN) objective}
\label{sec:gan}
The runs so far used only MSE. Every recipe was retrained with the VQGAN objective --- an
LPIPS perceptual term plus a PatchGAN adversarial loss (Esser et al.\ 2021) --- wiring otherwise
unchanged. We call them \textbf{A+, B+, C+, D+}. Three questions follow, and this Part answers
them in order: does perceptual quality improve (\S\ref{sec:gan})? does the GAN, rather than the
prior, explain the mixture (\S\ref{sec:twobytwo})? does the \emph{codebook}, rather than the
prior, explain it (\S\ref{sec:vqcontrol})?

\subsection*{Reconstruction quality: FID and LPIPS collapse dramatically}
\begin{table}[h]
\centering
\caption{Full-model reconstruction on the \textbf{full val set ($3925$ images)}, same protocol
(batch-wise, torchmetrics rFID/KID, AlexNet LPIPS) for all rows including the VAE, VQ-VAE, RAE and
VQGAN baselines. The GAN/LPIPS objective is a \textbf{perceptual} win: \textbf{rFID drops
$\sim$6--8$\times$ and LPIPS $\sim$4$\times$} vs the MSE runs. It also \textbf{flattens the whole
$2\times2$}: A+, B+, C+ and D+ land within $0.2$ rFID of each other, so neither $\beta$ nor the
prior costs anything in reconstruction. \textbf{That flatness has a ceiling}: \textbf{E+}
($\beta{=}1.0$, one decade past B+) finally \emph{does} pay, at rFID $4.94$ and LPIPS $0.066$
(\S\ref{sec:betaceiling}). The VAE/VQ-VAE rows form a complete \emph{single-branch
ablation} (all matched: 8-ch latent, LPIPS+GAN, $200$ ep): \textbf{VAE} $=$ continuous branch only,
\textbf{VQ-VAE} $=$ discrete branch only. \textbf{RAE} and \textbf{VQGAN} are external-recipe
models (their own encoder/latent designs), each \emph{retrained on Imagenette for $200$ epochs}
--- reference points rather than the matched single-branch ablation.}
\label{tab:recon}
\begin{tabular}{lccc ccc}
\toprule
 & \multicolumn{3}{c}{fidelity} & \multicolumn{3}{c}{\textbf{perceptual}}\\
\cmidrule(r){2-4}\cmidrule(l){5-7}
Run & MSE\,$\downarrow$ & PSNR\,$\uparrow$ & SSIM\,$\uparrow$ & LPIPS\,$\downarrow$ & \textbf{rFID}\,$\downarrow$ & KID\,$\downarrow$\\
\midrule
A \ ($\beta{=}.001$, std, MSE) & $0.0013$ & \textbf{$29.37$} & \textbf{$0.878$} & $0.112$ & $24.4$ & $0.0099$\\
B \ ($\beta{=}.1$, std, MSE)   & $0.0020$ & $27.29$ & $0.814$ & $0.189$ & $40.0$ & $0.0194$\\
C \ ($\beta{=}.1$, code, MSE)  & $0.0017$ & $27.98$ & $0.835$ & $0.172$ & $31.9$ & $0.0143$\\
D \ ($\beta{=}.001$, code, MSE)& $0.0015$ & $28.89$ & $0.869$ & $0.117$ & $22.7$ & $0.0090$\\
\midrule
\textbf{A+} (std, $\beta{=}.001$) & $0.0020$ & $27.26$ & $0.834$ & \good{$0.043$} & \good{$4.02$} & \good{$0.0001$}\\
\textbf{B+} (std, $\beta{=}.1$)   & $0.0022$ & $26.94$ & $0.829$ & \good{$0.043$} & \good{$3.70$} & \good{$0.0000$}\\
\textbf{C+} (code, $\beta{=}.1$)  & $0.0021$ & $27.34$ & $0.835$ & \good{$0.043$} & \good{$3.86$} & \good{$0.0001$}\\
\textbf{D+} (code, $\beta{=}.001$)& $0.0021$ & $27.27$ & $0.833$ & \good{$0.044$} & \good{$3.97$} & \good{$0.0001$}\\
\midrule
E+ (std, $\beta{=}1.0$)           & $0.0032$ & $25.23$ & $0.783$ & $0.066$ & \bad{$4.94$} & $0.0001$\\
\midrule
VAE (continuous only, $+$LPIPS/GAN) & $0.0021$ & $27.27$ & $0.832$ & \good{$0.043$} & \good{$4.08$} & $0.0001$\\
VQ-VAE (codes only, $+$LPIPS/GAN) & $0.0034$ & $25.07$ & $0.776$ & $0.074$ & $5.93$ & $0.0005$\\
RAE (SOTA baseline, $200$ ep) & $0.0197$ & $17.18$ & $0.461$ & $0.210$ & $7.64$ & $0.0004$\\
VQGAN (Esser et al., paper recipe, retrained) & $0.0131$ & $19.02$ & $0.517$ & $0.241$ & $24.59$ & $0.0057$\\
\bottomrule
\end{tabular}
\end{table}

\textbf{Plain-words result.} The adversarial + perceptual objective does exactly what it is meant
to: \textbf{rFID $24$--$40 \to \approx4$} and \textbf{LPIPS $0.11$--$0.19 \to 0.043$} (KID
$\to 0.0001$), i.e.\ far more realistic reconstructions (Fig.~\ref{fig:reconCg}), while PSNR/SSIM
barely move --- the GAN optimizes realism, not pixel fidelity. A second, subtler finding: the
\emph{whole} $2\times2$ collapses onto one reconstruction level. Without the GAN, D beat C and A
beat B by a wide margin; with it, A+, B+, C+ and D+ span a mere $0.32$ rFID
($3.70$--$4.02$, all $\approx$$0.043$ LPIPS). Mechanistically the GAN drives every run's
continuous branch to high rate (posterior std $\to 0.006$--$0.1$, near-deterministic), so the
choice of $\beta$ --- and, for reconstruction, the choice of prior --- stops mattering. This
confirms the caveat from \S\ref{sec:codeprior}: judge these models on FID/LPIPS, not PSNR/SSIM.
\textbf{It also sets up the key experiment:} if all four reconstruct identically, any difference in
latent structure is attributable to the prior alone. That is \S\ref{sec:twobytwo}.

\textbf{A complete single-branch ablation (VAE vs VQ-VAE).} Two matched baselines --- both with the
\emph{same} 8-ch latent, LPIPS and GAN, $200$ epochs --- isolate each branch. The \textbf{VAE}
(continuous branch only, \emph{no} codebook/prior) reaches rFID $4.08$; the \textbf{VQ-VAE} (discrete
branch only, \emph{no} continuous/prior) reaches rFID $5.93$. Reading the two together is the honest
story of where reconstruction lives:
\begin{itemize}
  \item \textbf{Removing the \emph{discrete} branch barely moves reconstruction:} dual D+ $3.97
  \to$ VAE $4.08$ (LPIPS $0.044\to0.043$, SSIM $0.833\to0.832$ --- a tie). \emph{A matched continuous
  VAE reconstructs as well as the dual model.} So the dual design is \textbf{not} a reconstruction win
  over a plain VAE.
  \item \textbf{Removing the \emph{continuous} branch clearly hurts:} dual D+ $3.97 \to$ VQ-VAE $5.93$
  ($+\!2$ rFID, LPIPS $0.044\to0.074$; Fig.~\ref{fig:vqrecon}). The continuous branch carries the fine
  detail; codes alone are softer.
\end{itemize}
So the reconstruction is carried by the \emph{continuous} branch, and the dual model keeps that
VAE-level fidelity while \textbf{adding the discrete codes and the Gaussian-mixture latent at
essentially no rFID cost}. The honest headline is therefore not ``the dual reconstructs better'' but
\textbf{``the dual gets a structured, code-organized, generatively-usable latent for free, at
matched reconstruction''} --- its distinctive value is the codes/mixture
(\S\ref{sec:downstream},~\S\ref{sec:mcm}), not pixels.
Both single-branch baselines still beat RAE ($7.64$).

\begin{figure}[h]
\centering
\begin{subfigure}{0.49\textwidth}
  \includegraphics[width=\linewidth]{\runVAE/vae_reconstruction_grid.pdf}
  \caption{VAE (continuous only): rFID $4.08$, ties the dual model.}
\end{subfigure}\hfill
\begin{subfigure}{0.49\textwidth}
  \includegraphics[width=\linewidth]{\runVQ/vqvae_reconstruction_grid.pdf}
  \caption{VQ-VAE (codes only): rFID $5.93$, slightly softer.}
\end{subfigure}
\caption{Matched single-branch baselines on $5$ val images (original / reconstruction). The
continuous-only VAE is indistinguishable from the dual model; the codes-only VQ-VAE is softer ---
the continuous branch carries the fine detail.}
\label{fig:vqrecon}
\end{figure}

\textbf{Comparison to SOTA (RAE).} As an external reference we ran the identical full-val
protocol on \textbf{RAE} (Representation Autoencoders for diffusion transformers), trained on the
same data for the same $200$ epochs. Our perceptual models reconstruct \textbf{clearly better}:
rFID \textbf{$3.9$ vs $7.6$} ($\sim$$2\times$), LPIPS \textbf{$0.043$ vs $0.210$} ($\sim$$5\times$),
and PSNR/SSIM far higher ($27.3/0.83$ vs $17.2/0.46$). \emph{Fair-comparison note:} RAE builds its
latent from a \emph{frozen semantic encoder} chosen to make the latent good for \emph{diffusion
generation}, so it deliberately trades pixel fidelity for a generation-ready, semantic latent ---
which is why its PSNR/SSIM are low by design. The win here is therefore specifically on the
\textbf{reconstruction-quality axis}; on the \emph{semantic} axis a frozen-DINO latent like RAE's
would lead (cf.\ \S\ref{sec:downstream}, where our codes' probe accuracy is a modest
$\approx$$0.44$). So the honest headline is: \textbf{our AE reconstructs substantially better than
RAE at equal budget while also carrying a Gaussian-mixture latent} --- not that it is a better
representation for downstream.

\begin{figure}[h]
\centering
\includegraphics[width=0.62\textwidth]{\runCg/ablation_grid_\tagCg.pdf}
\caption{\textbf{C+} branch-ablation reconstructions. \texttt{full} and \texttt{vq\_only} are
crisp and near-identical to the originals (LPIPS $0.043$ / $0.088$); \texttt{vanilla\_only} is
featureless (the codes carry the image). The equivalent A+/B+/D+ figures are in their folders.}
\label{fig:reconCg}
\end{figure}

\subsection*{Does the mixture survive the GAN under the code prior? Yes --- sharper components,
smoother envelope}
\begin{table}[h]
\centering
\caption{Mixture diagnostics before vs after the GAN, code-centered prior. Per-code structure
\emph{strengthens} (separability, silhouette both up), but the aggregate marginal becomes
\emph{smoother} (the global multimodality tests soften).}
\label{tab:gmmgan}
\small
\begin{tabular}{lcccc}
\toprule
Metric & \textbf{C} & \textbf{C+} & \textbf{D} & \textbf{D+}\\
\midrule
Active units                     & \good{100\%} & \good{100\%} & \good{100\%} & \good{100\%}\\
Codes used                       & $256/256$ & $256/256$ & $256/256$ & $256/256$\\
\textbf{Separability ratio} ($\gg1$) & $2.72$ & \good{$3.04$} & $2.35$ & \good{$3.03$}\\
Silhouette by code ($>0$)        & $0.13$ & \good{$0.19$} & $0.16$ & \good{$0.18$}\\
Code-anchored vs.\ single LL ($>0$) & \good{$+1.27$} & \good{$+0.17$} & \good{$+0.63$} & \good{$+0.16$}\\
GMM vs.\ single BIC ($<0$)       & \good{$-10347$} & \bad{$+16503$} & \good{$-3508$} & \bad{$+16885$}\\
Hartigan dip $p$ ($<0.05$)       & \good{$0.0007$} & \bad{$0.99$} & \good{$0.005$} & \bad{$1.0$}\\
Posterior std / KL\,nats/elem    & $.22 / 1.0$ & $.05 / 2.6$ & $.05 / 2.6$ & $.04 / 2.7$\\
\bottomrule
\end{tabular}
\end{table}

\textbf{Plain-words result.} No collapse (all $256$ codes used, $100\%$ active units in both).
The \emph{per-code} structure actually \textbf{improves}: separability climbs to $\sim$$3.0$
(codeword spacing grows to $\sim$$1.5$ while the residual stays $\sim$$0.5$) and the silhouette
rises to $\sim$$0.18$, and the code-anchored mixture still beats a single Gaussian in held-out
likelihood. So the codewords remain distinct, well-separated anchors (Fig.~\ref{fig:gmmCg},
colored clusters). But the \emph{aggregate} multimodality tests soften: the dip test is no longer
significant and the free-fit GMM BIC now prefers a single Gaussian, matching the smooth grey
envelope filling the centre of Fig.~\ref{fig:gmmCg}. The reading: the adversarial objective keeps
the codes as sharp mixture components while smoothing the \emph{overall} latent envelope toward a
single Gaussian --- the mixture survives at the component level, not as a globally multimodal
marginal. (Branch ablation is unchanged in spirit: the codes carry the image, and the continuous
branch now contributes \emph{more perceptually} --- it roughly halves LPIPS, $0.088\to0.043$ ---
because the GAN's fine detail lives there.)

\begin{figure}[h]
\centering
\includegraphics[width=0.60\textwidth]{\runCg/gmm_pca_codes_\tagCg.pdf}
\caption{\textbf{C+} per-location latent (PCA 2-D), colored by code. Top codes form tight,
well-separated clusters on their centroids (separability $3.04$), but the aggregate (grey) is a
smooth central envelope --- so global multimodality tests read as unimodal even though the
components are sharp.}
\label{fig:gmmCg}
\end{figure}

\section{Completing the $2\times2$: the GAN is \emph{not} a substitute for the prior (A+/B+)}
\label{sec:twobytwo}
Up to here the perceptual objective had only ever been measured \emph{under the code-centered
prior} (C+/D+), which leaves the headline of \S\ref{sec:gan} confounded: is the surviving mixture
a property of the prior, or would adversarial pressure have organized the codes anyway? Runs
\textbf{A+} and \textbf{B+} close that hole. They are byte-for-byte the C+/D+ recipe --- EMA
codebook, residual-quantization depth $2$, k-means codebook init, LPIPS weight, PatchGAN schedule,
$200$ epochs --- with \emph{two flags flipped}: \texttt{residual\_continuous} and
\texttt{component\_prior} off, i.e.\ the standard $\mathcal{N}(0,I)$ prior of Runs~A/B.
That makes A+/B+/C+/D+ a clean $2\times2$ of \textbf{prior} $\times$ $\beta$ at fixed objective.

\begin{table}[h]
\centering
\caption{The $2\times2$ under the perceptual objective. \textbf{Reconstruction is flat across the
whole table} (rFID $3.70$--$4.02$, LPIPS $0.043$--$0.044$) --- so the prior costs nothing --- but
the \emph{mixture} appears in exactly the two code-prior cells. The decisive row is the
code-anchored held-out likelihood, which is the direct test of ``are the codewords the mixture
means?'': positive only with the code-centered prior.}
\label{tab:twobytwo}
\small
\begin{tabular}{lcccc}
\toprule
 & \multicolumn{2}{c}{standard $\mathcal{N}(0,I)$ prior} & \multicolumn{2}{c}{code-centered prior}\\
\cmidrule(r){2-3}\cmidrule(l){4-5}
Metric & \textbf{A+} ($\beta{=}.001$) & \textbf{B+} ($\beta{=}.1$) & \textbf{D+} ($\beta{=}.001$) & \textbf{C+} ($\beta{=}.1$)\\
\midrule
rFID $\downarrow$                & $4.02$ & $3.70$ & $3.97$ & $3.86$\\
LPIPS $\downarrow$               & $0.043$ & $0.043$ & $0.044$ & $0.043$\\
SSIM $\uparrow$                  & $0.834$ & $0.829$ & $0.833$ & $0.835$\\
\midrule
Codes used                       & \bad{$101/256$} & $190/256$ & \good{$256/256$} & \good{$256/256$}\\
Codeword spacing                 & $0.58$ & $1.33$ & $1.52$ & $1.49$\\
Residual std                     & \bad{$2.40$} & $1.11$ & \good{$0.50$} & \good{$0.49$}\\
\textbf{Separability ratio} ($\gg1$) & \bad{$0.24$} & \bad{$1.19$} & \good{$3.03$} & \good{$3.04$}\\
Silhouette by code ($>0$)        & \bad{$-0.06$} & $0.04$ & \good{$0.18$} & \good{$0.19$}\\
\textbf{Code-anchored vs.\ single LL} ($>0$) & \bad{$-3.12$} & \bad{$-0.18$} & \good{$+0.16$} & \good{$+0.17$}\\
Effective \# components          & $87$ & $147$ & $226$ & $226$\\
\bottomrule
\end{tabular}
\end{table}

\textbf{Plain-words result --- three findings.}
\begin{itemize}
  \item \textbf{The GAN erases every reconstruction difference, including the prior's.} All four
  cells land at rFID $\approx$$4$ and LPIPS $\approx$$0.043$. So the code-centered prior is
  \emph{free}: it neither helps nor hurts pixels once the objective is perceptual. This is a
  stronger version of the ``$\beta$ gap disappears'' claim --- the \emph{prior} gap disappears too,
  on the reconstruction axis only.
  \item \textbf{The mixture does not appear without the prior.} Adversarial pressure does not
  organize the codes into components: A+ is still one blob (separability $0.24$, silhouette
  $-0.06$, Fig.~\ref{fig:gmmAg}a) and B+ only reaches the same ``emerging'' state as B
  ($1.19$, Fig.~\ref{fig:gmmAg}b). Crucially, the code-anchored held-out likelihood --- the one
  test that asks Claim~B directly --- stays \textbf{negative} in both ($-3.12$, $-0.18$) and is
  positive in both code-prior cells. \emph{The mixture is a property of the prior's shape, full
  stop.}
  \item \textbf{Without the prior, the GAN actively \emph{hurts} the codebook.} A+ uses only
  $101$ of $256$ codewords (Run~A used $194$) with an effective component count of $87$, and its
  continuous residual balloons to $2.40$ --- ten times the code spacing. Freed by the weak KL and
  pushed by the adversarial detail term, the continuous branch absorbs everything and the codebook
  half-dies. Under the code-centered prior the same objective leaves all $256$ codes alive with
  $226$ effective components.
\end{itemize}
Taken with \S\ref{sec:vqcontrol} (a codebook alone does not produce a mixture either), this rules
out both alternative explanations: neither the adversarial objective nor the codebook makes the
latent a mixture --- only the code-centered prior does, and it does so at no reconstruction cost.

\begin{figure}[h]
\centering
\begin{subfigure}{0.49\textwidth}
  \includegraphics[width=\linewidth]{\runAg/gmm_pca_codes_\tagAg.pdf}
  \caption{\textbf{A+} (std prior, $\beta{=}.001$, GAN): centroids piled in the middle of one
  large cloud; separability $0.24$, only $101$ codes alive.}
\end{subfigure}\hfill
\begin{subfigure}{0.49\textwidth}
  \includegraphics[width=\linewidth]{\runBg/gmm_pca_codes_\tagBg.pdf}
  \caption{\textbf{B+} (std prior, $\beta{=}.1$, GAN): centers spread, points loosely organize
  --- but still one envelope; separability $1.19$, anchored LL $-0.18$.}
\end{subfigure}
\caption{The standard prior \emph{with} the perceptual objective. Compare the tight, separated,
likelihood-confirmed components of C+ (Fig.~\ref{fig:gmmCg}) trained with the \emph{same} objective
and the \emph{same} codebook machinery. The only difference is where the Gaussian is centered.}
\label{fig:gmmAg}
\end{figure}

\section{The $\beta$ ceiling: how far the standard prior can be pushed (E+)}
\label{sec:betaceiling}
\S\ref{sec:twobytwo} showed the standard prior does not produce a mixture at $\beta\in\{.001,.1\}$.
The obvious objection is that we simply did not push $\beta$ hard enough: B+ was already
\emph{moving} in the right direction ($0.24\to1.19$ separability from A+), so perhaps one more
decade closes the gap and the code-centered prior is unnecessary. Run~\textbf{E+} tests exactly
that. It is the B+ recipe with a single number changed --- $\beta=1.0$, ten times B+ and a
thousand times A+ --- keeping the standard $\mathcal{N}(0,I)$ prior, the EMA codebook, RQ depth
$2$, k-means init, LPIPS weight, PatchGAN schedule and $200$ epochs. It extends the standard-prior
$\beta$ ladder to three points: $.001\to.1\to1.0$.

\begin{table}[h]
\centering
\caption{The standard-prior $\beta$ ladder (A+, B+, E+) against the code-centered cells (C+, D+).
Every \emph{geometric} proxy improves monotonically with $\beta$ and E+ nearly reaches the
code-prior values --- but the \textbf{decisive likelihood test never changes sign}, and E+ pays for
the approach with reconstruction and with the first posterior collapse in this report.}
\label{tab:betaceiling}
\small
\begin{tabular}{lccc cc}
\toprule
 & \multicolumn{3}{c}{standard $\mathcal{N}(0,I)$ prior --- $\beta$ ladder} & \multicolumn{2}{c}{code-centered}\\
\cmidrule(r){2-4}\cmidrule(l){5-6}
Metric & \textbf{A+} ($\beta{=}.001$) & \textbf{B+} ($\beta{=}.1$) & \textbf{E+} ($\beta{=}1.0$) & \textbf{D+} ($\beta{=}.001$) & \textbf{C+} ($\beta{=}.1$)\\
\midrule
rFID $\downarrow$              & \good{$4.02$} & \good{$3.70$} & \bad{$4.94$} & \good{$3.97$} & \good{$3.86$}\\
LPIPS $\downarrow$             & \good{$0.043$} & \good{$0.043$} & \bad{$0.066$} & \good{$0.044$} & \good{$0.043$}\\
SSIM $\uparrow$                & $0.834$ & $0.829$ & \bad{$0.783$} & $0.833$ & $0.835$\\
\midrule
Active units ($100\%$ wanted)  & \good{$100\%$} & \good{$100\%$} & \bad{$87.5\%$} & \good{$100\%$} & \good{$100\%$}\\
Codes used                     & \bad{$101/256$} & $190/256$ & \good{$256/256$} & \good{$256/256$} & \good{$256/256$}\\
Codeword spacing               & $0.58$ & $1.33$ & $4.11$ & $1.52$ & $1.49$\\
Residual std                   & \bad{$2.40$} & $1.11$ & $1.50$ & \good{$0.50$} & \good{$0.49$}\\
Effective \# components        & $87$ & $147$ & $207$ & $226$ & $226$\\
\midrule
\textbf{Separability ratio} ($\gg1$) & \bad{$0.24$} & \bad{$1.19$} & $2.74$ & \good{$3.03$} & \good{$3.04$}\\
Silhouette by code ($>0$)      & \bad{$-0.06$} & $0.04$ & $0.17$ & \good{$0.18$} & \good{$0.19$}\\
\textbf{Code-anchored vs.\ single LL} ($>0$) & \bad{$-3.12$} & \bad{$-0.18$} & \bad{$-0.08$} & \good{$+0.16$} & \good{$+0.17$}\\
Hartigan dip $p$               & --- & --- & $0.95$ & $1.0$ & $0.99$\\
\bottomrule
\end{tabular}
\end{table}

\textbf{Plain-words result --- the standard prior gets close, and still fails where it counts.}
\begin{itemize}
  \item \textbf{Every geometric proxy improves, monotonically, and nearly closes the gap.}
  Along $\beta=.001\to.1\to1.0$: separability $0.24\to1.19\to\mathbf{2.74}$, silhouette
  $-0.06\to0.04\to\mathbf{0.17}$, live codes $101\to190\to\mathbf{256}$, effective components
  $87\to147\to\mathbf{207}$. E+ is within $10\%$ of C+/D+ on all four. Read only this block and you
  would conclude that enough KL pressure does replace the code-centered prior.
  \item \textbf{But the decisive test does not follow.} The code-anchored held-out likelihood ---
  fix the mixture means \emph{at the codewords} and ask whether that beats a single Gaussian ---
  goes $-3.12\to-0.18\to\mathbf{-0.08}$. It approaches zero and \textbf{stops there}: still
  negative, still on the wrong side, after a $1000\times$ sweep of the only knob available. Both
  code-prior cells are positive ($+0.16$, $+0.17$) at \emph{either} $\beta$. The dip test
  ($p=0.95$) and BIC ($-21701$, favouring a single Gaussian) agree the marginal stays one blob.
  \item \textbf{The improvement is the VQ-VAE artifact, not a mixture.} E+'s geometry is
  diagnostic: codeword spacing balloons to $4.11$ (B+: $1.33$) while the residual also grows to
  $1.50$. That is precisely the plain VQ-VAE's signature (\S\ref{sec:vqcontrol}: spacing $4.42$,
  residual $1.31$, separability $3.37$, anchored LL $-0.97$) --- a large latent finely tiled by
  256 codewords, which inflates spacing-over-residual ratios without making the codewords into
  mixture \emph{means}. Fig.~\ref{fig:gmmEg} shows it directly: the coloured per-code clusters do
  separate along PC1, but the grey aggregate is one broad, gap-free cloud.
  \item \textbf{And it is the first run in this report where Claim~A fails.}
  Active units fall to $\mathbf{87.5\%}$: one of the eight latent channels is completely dead
  (per-channel KL $0.0009$ nats, active fraction $0.0$), while posterior std rises to $0.673$ ---
  an order of magnitude above the other GAN runs' $0.04$--$0.05$. At $\beta=1.0$ the KL term
  finally wins the tug-of-war and starts switching the continuous branch off, exactly the
  failure mode \S\ref{sec:collapse} found absent everywhere else.
  \item \textbf{It also costs the reconstruction the GAN had made free.} rFID $3.70\to4.94$ and
  LPIPS $0.043\to0.066$ --- E+ is the worst reconstructor of every LPIPS/GAN model in
  Table~\ref{tab:recon}, worse than the single-branch VAE ($4.08$) and the single-branch
  VQ-VAE is its only company ($5.93$). The image-space signature agrees: E+ has the \emph{lowest}
  cluster coherence in the whole grid ($22.4\%$, Table~\ref{tab:seg}) despite spending the most
  codewords per image ($212.5$) --- a larger vocabulary describing a more fragmented picture.
\end{itemize}

\textbf{Why this matters for the report's headline.} \S\ref{sec:mixture} established a
rate/separability \emph{trade-off} for the standard prior under the MSE objective, and
\S\ref{sec:gan} showed the GAN hides it (all four $2\times2$ cells reconstruct identically).
E+ shows the trade-off was never removed --- only pushed out of the measured range. Extend $\beta$
one decade and it reappears in full: you buy separability $1.19\to2.74$ and pay $1.24$ rFID, a dead
channel, and a $34\%$ worse LPIPS, \emph{and you still do not get a code-anchored mixture}. The
code-centered prior reaches strictly better structure (separability $3.03$, anchored LL $+0.16$)
at $\beta=0.001$, i.e.\ at a thousandth of the KL pressure, with rFID $3.97$ and $100\%$ active
units. \textbf{The two mechanisms are not substitutes: $\beta$ moves the latent's scale, the prior
moves its shape, and only shape produces a mixture.}

\begin{figure}[h]
\centering
\includegraphics[width=0.60\textwidth]{\runEg/gmm_pca_codes_\tagEg.pdf}
\caption{\textbf{E+} (standard prior, $\beta{=}1.0$, GAN) per-location latent (PCA 2-D), colored by
code, $\bigstar=$ code centroids. The strongest standard-prior geometry in the report --- centroids
genuinely spread along PC1, separability $2.74$ --- and yet the grey aggregate is a single broad
cloud with no gaps, the dip test reads unimodal ($p=0.95$), and the codeword-anchored mixture is
still \emph{worse} than one Gaussian ($-0.08$). Compare the plain VQ-VAE (Fig.~\ref{fig:gmmVQ}),
which produces the same picture for the same reason, and C+ (Fig.~\ref{fig:gmmCg}), which does not.}
\label{fig:gmmEg}
\end{figure}

\section{Where the mixture comes from: the plain VQ-VAE control}
\label{sec:vqcontrol}
The VQ-VAE of Table~\ref{tab:recon} shares the codebook, RQ depth, LPIPS and GAN with D+ but
\emph{drops the continuous branch and the code-centered prior}. It is therefore the other half of
the control: \S\ref{sec:twobytwo} showed the GAN is not the cause; this shows the codebook is not
either. We run the \emph{identical} Claim-B analysis on the VQ-VAE's pre-quantization
encoder output $z_e$ --- the exact analog of the dual model's ``codeword $+$ offset'' latent, since
$z_e=e_k+(z_e-e_k)$ is a codeword plus a quantization residual (200 val images, same protocol as
Table~\ref{tab:gmm}).

\begin{table}[h]
\centering
\caption{Gaussian-mixture diagnostics: plain VQ-VAE vs.\ the matched dual GAN models. The
\emph{geometric} proxies (separability, silhouette) even look \emph{better} for the VQ-VAE --- but
they are inflated by fine tiling of a large-scale latent (its spacing and residual are both
$\sim$$3\times$ larger). The one test that directly asks ``are the codewords the mixture
\emph{means}?'' --- the held-out likelihood of a codeword-anchored mixture vs.\ a single Gaussian
--- \textbf{flips sign}: positive for every code-prior dual run, \textbf{negative for the VQ-VAE}
(and, from Table~\ref{tab:twobytwo}, negative for A+/B+ too).}
\label{tab:vqcontrol}
\small
\begin{tabular}{lccc}
\toprule
Metric (want for a code-anchored GMM) & \textbf{VQ-VAE} & \textbf{C+} & \textbf{D+}\\
 & (codes only) & (dual, GAN) & (dual, GAN)\\
\midrule
Codes used                            & $256/256$ & $256/256$ & $256/256$\\
Codeword spacing                      & $4.42$ & $1.49$ & $1.52$\\
Residual std                          & $1.31$ & $0.49$ & $0.50$\\
Separability ratio ($\gg1$)           & $3.37$ & $3.04$ & $3.03$\\
Silhouette by code ($>0$)             & $0.25$ & $0.19$ & $0.18$\\
\textbf{Code-anchored vs.\ single LL} ($>0$) & \bad{$-0.97$} & \good{$+0.17$} & \good{$+0.16$}\\
GMM vs.\ single BIC ($<0$)            & \bad{$+4442$} & \bad{$+16503$} & \bad{$+16885$}\\
Hartigan dip $p$                      & $0.10$ & $0.99$ & $1.0$\\
\bottomrule
\end{tabular}
\end{table}

\textbf{Plain-words result.} Without the code-centered prior, the codebook is just a
$k$-means-style quantizer of a \emph{single} continuous Gaussian: 256 codewords tile one big blob
(Fig.~\ref{fig:gmmVQ}). That fine tiling makes the naive geometric scores look strong (each cell is
small relative to the spacing between cells), which is exactly why separability and silhouette are
\emph{not} sufficient evidence for a mixture --- they read high even here. The decisive, generative
test is the \textbf{code-anchored held-out likelihood}: fix the mixture means to the codewords and
ask whether that beats a single Gaussian. For every code-prior dual model it does ($+0.16$ to
$+1.27$; the prior $\mathcal{N}(e_k,\sigma_k^2)$ was \emph{built} to put the codeword at the mean),
but for the VQ-VAE it is \textbf{negative} ($-0.97$): a codeword-anchored mixture explains the
latent \emph{worse} than one Gaussian, because approximating a single Gaussian with 256
sub-Gaussians only adds parameters without adding structure. The dip test (unimodal, $p=0.10$) and
BIC (prefers a single Gaussian) agree the marginal is one blob. \textbf{So the Gaussian mixture is
produced by the code-centered prior, not by the codebook} --- a VQ-VAE with the same codebook does
not have it.

\begin{figure}[h]
\centering
\includegraphics[width=0.60\textwidth]{\runVQ/gmm_pca_honest_\tagVQ.pdf}
\caption{\textbf{Plain VQ-VAE} pre-quantization latent $z_e$ (PCA 2-D), colored by code,
$\bigstar=$ code centroids. The codes tile \emph{one} continuous Gaussian cloud (grey envelope
fills the centre, no gaps between components) --- the codeword-anchored mixture is worse than a
single Gaussian ($-0.97$), unlike the tight, likelihood-confirmed components of the code-prior dual
models (Figs.~\ref{fig:gmmC},~\ref{fig:gmmCg}).}
\label{fig:gmmVQ}
\end{figure}

\section{What the codes actually segment}
\label{sec:seg}
The mixture tests live in latent space. This section asks the same question \emph{on the image}:
if the codewords are genuine mixture components, do they land on coherent pieces of the picture?
We paint each of the $32\times32$ latent locations with the colour of the codeword it is assigned
to. Raw code colours are unreadable at $256$ categories, so we group the codebook into
$K{=}10$ super-clusters by $k$-means and colour those. Colours are per-model (the codebooks and
their groupings differ), so compare \emph{region structure within an image}, never colours across
columns. Two numbers accompany each panel: distinct codewords used in that image (lower $=$ less
fragmented) and \textbf{neighbour coherence} --- the fraction of adjacent patch pairs assigned to
the same super-cluster (higher $=$ larger, more contiguous regions).

\begin{figure}[h]
\centering
\includegraphics[width=\textwidth]{reports/code_cluster_seg_all_runs.pdf}
\caption{\textbf{Code-cluster segmentation across the prior$\times$objective grid.} Columns:
original, then Runs A, B (standard prior, MSE), \textbf{A+, B+, E+} (standard prior, LPIPS+GAN,
$\beta=.001/.1/1.0$) and C+, D+ (code-centered prior, LPIPS+GAN). Colour $=$ super-cluster
($k$-means, $K{=}10$) of the codeword assigned to each latent patch; per-column headers give the
mean over $60$ val images.
Every model finds the coarse figure/ground split (row~2: the parachute against the sky), and none
produces object-level segments --- a reconstruction codebook tiles \emph{appearance}, not objects.
\textbf{E+} is the image-space counterpart of \S\ref{sec:betaceiling}: the most codewords per image
of any run ($212.5$) and simultaneously the \emph{least} coherent map ($22.4\%$) --- maximal
KL pressure buys vocabulary, not structure (row~2, where its parachute region is visibly noisier
than C+/D+'s).
Regenerate with \texttt{tools/code\_cluster\_segmentation.py}; statistics in \texttt{reports/code\_cluster\_seg\_all\_runs.json}.}
\label{fig:seg}
\end{figure}

\begin{table}[h]
\centering
\caption{Code segmentation statistics over $60$ validation images ($K{=}10$ super-clusters).
``codes/img'' counts distinct codewords used per image; ``code coh.'' and ``cluster coh.'' are
neighbour agreement at codeword and super-cluster level.}
\label{tab:seg}
\small
\begin{tabular}{lccccccc}
\toprule
 & \textbf{A} & \textbf{B} & \textbf{A+} & \textbf{B+} & \textbf{E+} & \textbf{C+} & \textbf{D+}\\
 & (std,$.001$) & (std,$.1$) & (std,$.001$,GAN) & (std,$.1$,GAN) & (std,$1.0$,GAN) & (code,$.1$,GAN) & (code,$.001$,GAN)\\
\midrule
Distinct codes / image & $168.1$ & $176.6$ & \bad{$96.5$} & $150.2$ & \good{$212.5$} & \good{$212.3$} & \good{$209.8$}\\
Code-level coherence   & $3.2\%$ & \good{$11.8\%$} & $5.2\%$ & $5.3\%$ & $5.2\%$ & $7.2\%$ & $7.4\%$\\
Cluster-level coherence& \good{$35.2\%$} & $32.0\%$ & $26.4\%$ & $23.5\%$ & \bad{$22.4\%$} & $26.8\%$ & $28.0\%$\\
\bottomrule
\end{tabular}
\end{table}

\textbf{Plain-words result --- and an honest negative.}
\begin{itemize}
  \item \textbf{No run segments objects.} At $8\times$ downsampling with $256$ codes and a pure
  reconstruction objective, the assignment map is dominated by local appearance (colour, texture,
  brightness), so it is speckled everywhere. The figure/ground structure \emph{is} visible where
  the image has it (row~2: the parachute is a distinct region against a uniform sky in every
  column), and it fades on cluttered images (rows 1, 3, 4). This is the same ceiling
  \S\ref{sec:downstream} finds numerically: a reconstruction latent encodes appearance, not
  semantics.
  \item \textbf{Coherence tracks $\beta$, not the prior.} The most spatially coherent codes belong
  to \textbf{Run~B} (code-level $11.8\%$, cluster-level $32.0\%$) --- the run where the strong KL
  forces the codes to carry the coarse image (\S\ref{sec:branches}). The GAN runs are all in a
  narrow $23$--$28\%$ band. So spatial coherence is a \emph{different} axis from mixture
  separability, and the code-centered prior does not, by itself, buy larger segments.
  \item \textbf{But codebook \emph{usage} tracks the prior exactly.} The code-prior GAN runs spend
  $\approx$$210$ distinct codewords per image; \textbf{A+ spends only $96$} --- less than half, the
  image-space signature of the half-dead codebook found in latent space
  (\S\ref{sec:twobytwo}: $101/256$ codes alive, $87$ effective components). Note this is a
  \emph{narrower vocabulary}, not a cleaner one: A+'s cluster coherence ($26.4\%$) is no higher
  than the code-prior runs', so it is describing the same picture with fewer, blurrier words.
\end{itemize}

\section{A negative result: the wavelet detail branch (C++/D++)}
\label{sec:wavelet}
The \texttt{swd\_dualvae} model gained fine texture from a \emph{wavelet detail branch}: the
continuous branch encodes the PRE-compression high-frequency Haar subbands (LH/HL/HH), which the
strided encoder would otherwise average away. We ported that branch into the C+/D+ recipe ---
but kept the \emph{code-centered KL} regularizer (unlike \texttt{swd\_dualvae}, which regularizes
the wavelet with SWD + a per-band variance). \textbf{C++}/\textbf{D++} = C+/D+ with the continuous
branch replaced by the wavelet detail branch. \textbf{The result is a clean negative: it hurt
every metric.}

\begin{table}[h]
\centering
\caption{Wavelet vs.\ no-wavelet under the perceptual objective (full model, \textbf{full val,
$3925$ images}). The wavelet \emph{raises} rFID and LPIPS and lowers PSNR/SSIM --- it did not add
clean detail.}
\label{tab:wavelet}
\begin{tabular}{lcccc}
\toprule
Run & PSNR\,$\uparrow$ & SSIM\,$\uparrow$ & LPIPS\,$\downarrow$ & \textbf{rFID}\,$\downarrow$\\
\midrule
C+ \ (no wavelet, $\beta{=}.1$)   & \good{$27.34$} & \good{$0.835$} & \good{$0.043$} & \good{$3.86$}\\
D+ \ (no wavelet, $\beta{=}.001$) & \good{$27.27$} & \good{$0.833$} & \good{$0.044$} & \good{$3.97$}\\
C++ (wavelet, $\beta{=}.1$)       & \bad{$25.77$}  & \bad{$0.800$}  & \bad{$0.059$}  & \bad{$4.51$}\\
D++ (wavelet, $\beta{=}.001$)     & \bad{$25.88$}  & \bad{$0.805$}  & \bad{$0.054$}  & \bad{$4.39$}\\
\bottomrule
\end{tabular}
\end{table}

\textbf{And it damages the mixture} (Table~\ref{tab:waveletgmm}). The wavelet $\Delta$ is
high-variance (residual std $1.6$ at $\beta{=}0.1$, \textbf{$3.15$} at $\beta{=}0.001$, vs
$\approx0.5$ for C+/D+), because a single per-code isotropic $\sigma_k$ cannot shape three
frequency bands of very different scale. At low $\beta$ (D++) the wavelet branch runs unconstrained
(KL $8.2$ nats, $|\mu|{=}2.4$) and its magnitude \emph{swamps the codeword spacing}
($3.15>2.61$), so \textbf{separability falls to $0.83<1$} and the silhouette goes negative --- the
GMM is largely lost (Fig.~\ref{fig:gmmDpp}, one smeared cloud). \emph{This is the same failure mode
as A+} (\S\ref{sec:twobytwo}): whenever the continuous branch's spread outgrows the codeword
spacing, the mixture dies --- exactly the condition \eqref{eq:sepc} predicts.

\begin{table}[h]
\centering
\caption{Wavelet weakens the Gaussian mixture; at low $\beta$ (D++) it collapses it.}
\label{tab:waveletgmm}
\small
\begin{tabular}{lcccc}
\toprule
Metric & \textbf{C+} & \textbf{C++} & \textbf{D+} & \textbf{D++}\\
\midrule
Separability ratio ($\gg1$)      & \good{$3.04$} & $1.98$ & \good{$3.03$} & \bad{$0.83$}\\
Silhouette by code ($>0$)        & \good{$0.19$} & $0.14$ & \good{$0.18$} & \bad{$-0.04$}\\
Code-anchored vs.\ single LL ($>0$) & \good{$+0.17$} & \bad{$-0.03$} & \good{$+0.16$} & \bad{$-1.08$}\\
Continuous residual std          & $0.49$ & $1.60$ & $0.50$ & $3.15$\\
KL nats / element                & $2.6$  & $1.8$  & $2.7$  & $8.2$\\
\bottomrule
\end{tabular}
\end{table}

\textbf{Why (reasoning).} The wavelet's benefit in \texttt{swd\_dualvae} was \emph{coupled to its
SWD + per-band-variance regularizer}, not the ``feed pre-compression detail'' idea alone. Under
the mismatched per-code KL, the branch pumps poorly-shaped high-frequency energy into the latent
--- which the decoder renders as texture the InceptionV3/LPIPS statistics read as \emph{artifacts},
not clean detail, so rFID/LPIPS get worse. The branch ablation confirms it is very load-bearing at
low $\beta$ (D++: continuous adds $+13$ PSNR; wavelet-only $0.53$ SSIM $>$ codes-only $0.36$),
i.e.\ the wavelet dominates and de-anchors the codes. \textbf{Takeaway:} to make the wavelet help
here it needs a \emph{band-aware} regularizer (SWD or per-band $\sigma$) --- the SW\_DUALVAE path
--- not the per-code GMM KL of the DUALVAE port. As it stands, \textbf{C+/D+ (no wavelet) remain
the best FID configs}.

\begin{figure}[h]
\centering
\includegraphics[width=0.58\textwidth]{\runDpp/gmm_pca_codes_\tagDpp.pdf}
\caption{\textbf{D++} (wavelet, $\beta{=}0.001$) per-location latent, PCA 2-D. The high-variance
wavelet $\Delta$ swamps the codeword spacing (separability $0.83<1$): centroids pile near the centre
and the codes smear into one cloud --- the clean mixture of D+ (Fig.~\ref{fig:gmmC}) is lost.}
\label{fig:gmmDpp}
\end{figure}

\part{Is the organized latent worth anything?}

\section{Downstream: frozen-feature probing and retrieval}
\label{sec:downstream}
A prettier latent is only interesting if it \emph{helps}. So we ask a concrete, honest
question: does the Gaussian-mixture organization make the features more useful for
\textbf{classification} and \textbf{image retrieval}? We test it with the two standard
``frozen-feature'' protocols, \textbf{evaluated on the validation set only}.

\subsection*{What the two protocols mean (in plain words)}
\begin{itemize}
  \item \textbf{Linear probe.} Freeze the model, turn each image into one feature vector, and
  train \emph{only} a simple linear classifier on top to predict the $10$ Imagenette classes.
  If the classes are easy to separate with a straight cut, the features are ``semantically
  organized''. We split the val set in two (stratified): fit the linear classifier on one part,
  report top-$1$ accuracy on the held-out part. Chance $=0.107$.
  \item \textbf{Retrieval mAP.} Use each val image as a query and rank all other val images by
  cosine similarity; a good feature puts same-class images first. We report mean Average
  Precision (mAP) over all queries (leave-one-out). Chance $=0.100$.
\end{itemize}
No part of the network is fine-tuned --- only the frozen features are read out, which is the
honest way to ask ``what did the model already learn?''.

\subsection*{Four features per model}
The latent is an $8$-channel $32\times32$ grid, so we pool it to one vector per image four ways:
\textbf{z\_sum} (mean-pooled full latent, $8$-d), \textbf{z\_vq} (mean-pooled discrete branch,
$8$-d), \textbf{z\_cont} (mean-pooled continuous branch, $8$-d), and \textbf{bagcodes} --- the
normalized histogram of which of the $256$ codewords each image uses ($256$-d, a learned
``bag-of-visual-words''). A matched standalone VAE (mean-pooled $8$-d latent, same LPIPS+GAN recipe)
is the baseline; it has no discrete features by construction.

\begin{table}[h]
\centering
\caption{\textbf{Linear-probe top-1 accuracy} (left) and \textbf{retrieval mAP} (right),
full val, frozen features, identical stratified split. Bold $=$ best feature per row.}
\label{tab:downstream}
\small
\setlength{\tabcolsep}{4pt}
\begin{tabular}{lcccc@{\hskip 1.5em}cccc}
\toprule
 & \multicolumn{4}{c}{Linear probe (acc, chance $0.107$)} & \multicolumn{4}{c}{Retrieval (mAP, chance $0.100$)}\\
\cmidrule(r){2-5}\cmidrule(l){6-9}
Run & z\_sum & z\_vq & z\_cont & bagcodes & z\_sum & z\_vq & z\_cont & bagcodes\\
\midrule
A \ (std, $\beta{=}.001$, MSE)   & $0.280$ & $0.239$ & $0.278$ & \textbf{$0.290$} & $0.137$ & $0.119$ & $0.137$ & \textbf{$0.151$}\\
B \ (std, $\beta{=}.1$, MSE)     & $0.273$ & $0.273$ & $0.234$ & \textbf{$0.317$} & $0.136$ & $0.135$ & $0.123$ & \textbf{$0.142$}\\
C \ (code, $\beta{=}.1$, MSE)    & $0.297$ & $0.303$ & $0.145$ & \textbf{$0.348$} & $0.129$ & $0.129$ & $0.105$ & \textbf{$0.144$}\\
D \ (code, $\beta{=}.001$, MSE)  & $0.285$ & $0.306$ & $0.166$ & \textbf{$0.365$} & $0.128$ & $0.129$ & $0.113$ & \textbf{$0.142$}\\
\midrule
\textbf{A+} (std, $\beta{=}.001$)  & $0.361$ & \bad{$0.247$} & $0.363$ & \textbf{$0.385$} & $0.151$ & \bad{$0.122$} & $0.151$ & \textbf{$0.149$}\\
\textbf{B+} (std, $\beta{=}.1$)    & $0.358$ & \bad{$0.292$} & $0.359$ & \textbf{$0.324$} & $0.151$ & $0.135$ & \textbf{$0.158$} & $0.153$\\
\textbf{C+} (code, $\beta{=}.1$)   & $0.360$ & $0.357$ & $0.298$ & \textbf{$0.417$} & $0.156$ & $0.155$ & $0.138$ & \textbf{$0.171$}\\
\textbf{D+} (code, $\beta{=}.001$) & $0.354$ & $0.357$ & $0.287$ & \good{\textbf{$0.438$}} & $0.157$ & $0.156$ & $0.136$ & \textbf{$0.170$}\\
\midrule
D+MCM (code, $+$masked codes) & $0.355$ & $0.361$ & $0.299$ & \good{\textbf{$0.441$}} & $0.157$ & $0.156$ & $0.136$ & \textbf{$0.171$}\\
VAE (matched, GAN)            & $0.358$ & --      & --      & --      & $0.155$ & --      & --      & --\\
\bottomrule
\end{tabular}
\end{table}

\begin{figure}[h]
\centering
\includegraphics[width=0.92\textwidth]{reports/downstream_gan/downstream_probe_acc.pdf}
\caption{Linear-probe accuracy by feature and model across every pipeline.
(\texttt{Agan}=A+, \texttt{Bgan}=B+, \texttt{Cgan}=C+, \texttt{Dgan}=D+, \texttt{Dmcm}=D+MCM.)
Dashed line $=$ chance.}
\label{fig:probe}
\end{figure}

\subsection*{What we see, and whether it is ``worth it''}
\begin{itemize}
  \item \textbf{Everything is well above chance but modest.} Probe accuracy $0.24$--$0.44$ vs
  chance $0.107$ ($2$--$4\times$); retrieval barely clears chance ($0.11$--$0.17$ vs $0.10$).
  This is the expected ceiling: a \emph{reconstruction} latent captures appearance, not
  semantics, and will not rival a self-supervised model --- as anticipated, and as
  Fig.~\ref{fig:seg} shows visually.
  \item \textbf{Two levers, and the GAN is the bigger one.} On the best feature
  (\texttt{bagcodes}), the GMM organization lifts A$\to$D ($0.290\!\to\!0.365$, $+26\%$) and the
  perceptual objective lifts D$\to$D+ ($0.365\!\to\!0.438$, $+20\%$) --- the single biggest step.
  Forcing \emph{realistic} reconstructions pressures the codebook into more object-coherent,
  more linearly-separable codes.
  \item \textbf{The two levers are complementary, and the prior's contribution survives the GAN.}
  This is what the new A+/B+ runs settle. Under the GAN, the code-centered prior is still worth
  $+0.05$ on \texttt{bagcodes} (A+ $0.385 \to$ D+ $0.438$) and $+0.11$ on \texttt{z\_vq}
  ($0.247 \to 0.357$) --- so the \emph{downstream} value of the prior is not an artifact of the
  MSE-only comparison in the A--D block.
  \item \textbf{The advantage is specifically in the codes, not the pooled latent.} Every GAN
  model --- both priors, both $\beta$, and the matched plain VAE --- scores
  \texttt{z\_sum} $0.354$--$0.361$ and mAP $\approx0.155$. The pooled continuous latent carries the
  same modest class signal with or without a codebook. The dual/GMM design pulls ahead only through
  the discrete \texttt{bagcodes}/\texttt{z\_vq} features \emph{the VAE does not have}
  ($0.438$ vs the VAE's $0.358$).
  \item \textbf{The continuous branch becomes \emph{less} useful under the code prior}
  (\texttt{z\_cont} $0.363\!\to\!0.287$ from A+ to D+): exactly as intended --- it carries only the
  small within-component detail, so little class signal lives there. Under the standard prior it
  carries everything, and \texttt{z\_cont} $\approx$ \texttt{z\_sum}.
  \item \textbf{One anomaly worth flagging.} \textbf{B+} breaks the monotone pattern: its
  \texttt{bagcodes} probe ($0.324$) is \emph{below} A+'s ($0.385$) despite using nearly twice as
  many codewords. Under the standard prior the code histogram is not a stable descriptor --- what
  the codes encode is not pinned to anything --- so its downstream quality is erratic in a way the
  code-prior runs are not (C+ $0.417$, D+ $0.438$, D+MCM $0.441$ are tightly ordered).
  \item \textbf{Retrieval barely moves.} mAP goes $0.14 \to 0.17$ from the MSE to the GAN runs and
  is flat within each block --- the latent is simply not semantic enough for retrieval, whatever
  the prior.
\end{itemize}
\textbf{Verdict.} \emph{Worth it, in a specific and honest sense.} The Gaussian-mixture latent
is not just prettier: it makes the bag-of-codes a measurably better linear-probe descriptor, both
without the GAN ($+26\%$, A$\to$D) and with it ($+14\%$, A+$\to$D+), and D+ --- already a
best-in-class reconstructor --- is also the best on the code features. But the gains are confined
to the discrete representation and to linear probing; retrieval is nearly unmoved, the pooled latent
ties a plain VAE, and the absolute level confirms a reconstruction latent is no substitute for
self-supervised pre-training. If downstream semantics is the real target, the codebook is the useful
part to build on --- which is what \S\ref{sec:cvi} and \S\ref{sec:mcm} try.

\subsection*{And the plain VQ-VAE's codes?}
The codes-only control of \S\ref{sec:vqcontrol} lets us ask whether the continuous/GMM branch also
makes the \emph{codes} better downstream, or only prettier. Running the same protocol on it and on
the matched D+: the shared GAN already makes the codes useful --- consistent with the finding that
the GAN is the big lever --- but the dual model's codes are \emph{measurably better}
(Table~\ref{tab:vqdownstream}), while retrieval is tied. So the branch that turns the codes into
genuine mixture means also makes them a modestly better probe descriptor, on top of the GAN.

\begin{table}[h]
\centering
\caption{Downstream (frozen features, full val, chance acc $0.107$, mAP $0.100$): plain VQ-VAE vs.\
matched dual D+. The dual model's codes probe better; retrieval is unmoved (as everywhere).}
\label{tab:vqdownstream}
\small
\begin{tabular}{lcc}
\toprule
Feature & \textbf{VQ-VAE} & \textbf{D+} (dual)\\
\midrule
\texttt{z\_vq} \ probe acc      & $0.340$ & \good{$0.357$}\\
\texttt{bagcodes} probe acc     & $0.406$ & \good{$0.438$}\\
\texttt{bagcodes} retrieval mAP & $0.170$ & $0.171$\\
\bottomrule
\end{tabular}
\end{table}

\section{Pushing the codes toward semantics I: cross-view code-invariance (D+CVI)}
\label{sec:cvi}
The downstream probe caps out because a \emph{reconstruction} latent encodes
appearance, not semantics --- the same ceiling the segmentation picture shows
(\S\ref{sec:seg}). Run~\textbf{D+CVI} tries to fix that with a lightweight auxiliary head,
aiming to \textbf{raise linear-probe/retrieval without sacrificing reconstruction or the mixture}.

\subsection*{What was added for semantics (how it was trained)}
D+CVI is exactly Run~D+ (code-centered prior, LPIPS+GAN, $\beta{=}0.001$) with \emph{one extra
loss term} --- the reconstruction path, decoder, GAN and KL are all unchanged. Per training step:
\begin{itemize}
  \item The clean image is encoded and reconstructed as usual (so \textbf{FID is untouched}).
  \item A second, \emph{appearance-augmented} view of the same image (color jitter, blur,
  grayscale, moderate crop) is encoded (one extra encoder pass, under \texttt{no\_grad} --- cheap).
  \item For both views we form the image-level \textbf{soft bag-of-codes} --- a distribution over
  the $256$ codewords, averaged over locations. The \textbf{cross-view code-invariance (CVI)} loss
  pulls the clean distribution toward the augmented one (cross-entropy, augmented view as a
  stop-grad target): \emph{the codes a patch is assigned to must not change under color/blur/crop.}
\end{itemize}
\textbf{The idea:} linear probing rewards \emph{invariance}, reconstruction rewards
\emph{appearance}, so CVI injects invariance into the \emph{codes} (content) while the continuous
branch keeps absorbing appearance (style) --- a content/style split. Reconstruction still forces
the codes to be informative, so the aux term cannot collapse them; the CVI weight is ramped from
$0$ after the codebook settles. (Implementation in \texttt{experiments/train\_dualvae.py}; design
in \texttt{reports/semantic\_codes\_design.md}.)

\subsection*{Result: no sacrifice --- but also no semantic gain}
\begin{table}[h]
\centering
\caption{D+ vs.\ D+CVI. CVI perfectly \emph{preserves} reconstruction and the mixture (as
designed), but does \emph{not} improve the downstream metrics it targeted.}
\label{tab:cvi}
\small
\begin{tabular}{lccc@{\hskip 1.5em}cc@{\hskip 1.5em}cc}
\toprule
 & \multicolumn{3}{c}{reconstruction} & \multicolumn{2}{c}{mixture} & \multicolumn{2}{c}{downstream (bag-of-codes)}\\
\cmidrule(r){2-4}\cmidrule(lr){5-6}\cmidrule(l){7-8}
Run & LPIPS$\downarrow$ & rFID$\downarrow$ & SSIM$\uparrow$ & separab.$\uparrow$ & silh.$\uparrow$ & probe acc$\uparrow$ & mAP$\uparrow$\\
\midrule
D+     & $0.044$ & $3.97$ & $0.833$ & $3.03$ & $0.18$ & $0.438$ & $0.170$\\
D+CVI  & \good{$0.044$} & \good{$3.92$} & \good{$0.833$} & \good{$3.03$} & \good{$0.17$} & \bad{$0.437$} & \bad{$0.170$}\\
\bottomrule
\end{tabular}
\end{table}

\textbf{Plain-words result.} CVI did exactly half of what it promised: it \textbf{sacrificed
nothing} --- rFID $3.97\!\to\!3.92$, LPIPS/SSIM unchanged, separability $3.03\!\to\!3.03$, the GMM
fully intact --- but it \textbf{did not improve semantics}: bag-of-codes probe accuracy
$0.438\!\to\!0.437$ and retrieval mAP $0.170\!\to\!0.170$ are unchanged.

\textbf{Why CVI was a no-op (and how to fix it).} The most likely cause is a \emph{scale
mismatch}: the reconstruction loss is \texttt{sum}-reduced (hundreds--thousands per image), so a
CVI weight of $0.5$ against a per-image-mean CVI ($\approx\!5$) is negligible --- the aux gradient
barely nudged the codes over $200$ epochs. It is not disproven, just \textbf{under-weighted}: a
scale-matched weight ($\sim$$10^2$--$10^3\times$ larger, or a mean-reduced reconstruction),
stronger augmentations, and the symmetric loss would give it a real chance. When pushed harder it
destabilized the codebook (perplexity drifting down, no plateau) without lifting the probe ---
which motivated the change of family in \S\ref{sec:mcm}.

\section{Pushing the codes toward semantics II: masked code modeling (D+MCM)}
\label{sec:mcm}
CVI belongs to the \emph{invariance} family of self-supervision: all its signal comes from
augmentation, and it needs an anti-collapse crutch (Sinkhorn/me-max) that \emph{fights} the codebook.
So we switched \emph{families} to \textbf{masked modeling} (the BEiT/iBOT recipe): predict the
code of a \emph{hidden} patch from its visible neighbours. It is augmentation-free, it \emph{uses} the
codebook as its tokenizer instead of policing it, and it has no collapse shortcut --- a constant or
coarse code prediction scores badly, so no Sinkhorn/me-max is needed.

\subsection*{What was added (how it was trained)}
\textbf{D+MCM} is Run~D+ with \emph{one extra auxiliary head}, the reconstruction/GAN path untouched.
Per step: the latent grid $z_{e}$ is formed as usual; a small ViT predictor sees the grid with
$\sim$$50\%$ of the $32\times32$ locations \textbf{masked} (replaced by a learned mask token) and, from
the visible ones, predicts a distribution over the $256$ codes at each masked location. The target is
the model's \emph{own} soft code assignment (the GMM responsibilities) at those locations, detached.
Cross-entropy at the masked locations flows back through the \emph{visible} tokens into the encoder,
pressuring the codes to be \textbf{predictable from context} (= object-part / content structure). The
weight is bounded (a ratio-anchor capped so the effective weight cannot run away --- the lesson from
CVI). Implementation in \texttt{models/masked\_predictor.py} + \texttt{experiments/train\_dualvae.py}.

\subsection*{Result: preserves everything, small honest gain, and it does not collapse}
\begin{table}[h]
\centering
\caption{D+ vs.\ D+MCM (all full val). MCM \emph{perfectly preserves} reconstruction and the mixture
(aux head only) and \textbf{does not collapse the codebook} (unlike CVI), while giving a small, real
gain on the discrete code features. The matched VAE's \texttt{z\_sum} is shown to make the point that
the pooled latent is tied across all three --- the dual edge is the codes.}
\label{tab:mcm}
\small
\begin{tabular}{lcc@{\hskip 1.2em}cc@{\hskip 1.2em}ccc}
\toprule
 & \multicolumn{2}{c}{reconstruction} & \multicolumn{2}{c}{mixture} & \multicolumn{3}{c}{downstream probe acc$\uparrow$}\\
\cmidrule(r){2-3}\cmidrule(lr){4-5}\cmidrule(l){6-8}
Run & rFID$\downarrow$ & LPIPS$\downarrow$ & separab.$\uparrow$ & silh.$\uparrow$ & z\_sum & z\_vq & bagcodes\\
\midrule
D+     & $3.97$ & $0.044$ & $3.03$ & $0.18$ & $0.354$ & $0.357$ & $0.438$\\
D+MCM  & \good{$3.85$} & \good{$0.043$} & \good{$2.99$} & $0.14$ & $0.355$ & \good{$0.361$} & \good{$0.441$}\\
\midrule
VAE (matched, GAN) & $4.08$ & $0.043$ & -- & -- & $0.358$ & -- & --\\
\bottomrule
\end{tabular}
\end{table}

\textbf{Plain-words result.} MCM did what it was designed to do and nothing it wasn't:
\begin{itemize}
  \item \textbf{No sacrifice.} rFID $3.97\!\to\!3.85$, LPIPS unchanged, the Gaussian mixture intact
  (separability $3.03\!\to\!2.99$, code-anchored held-out likelihood still $>$ single, $256/256$ codes,
  $100\%$ active units) --- exactly as a reconstruction-untouching aux head should.
  \item \textbf{It does \emph{not} collapse the codebook.} This is the concrete advantage over CVI: the
  monitored \texttt{MSM Accuracy} (fraction of hidden patches whose code the context predicts) rose
  steadily while codebook perplexity stayed healthy --- there is no ``ignore codes'' shortcut, so the
  objective never fights the vocabulary.
  \item \textbf{A small, real semantic gain --- on the codes.} Bag-of-codes probe
  $0.438\!\to\!0.441$, \texttt{z\_vq} $0.357\!\to\!0.361$, \texttt{z\_cont} $0.287\!\to\!0.299$;
  retrieval flat ($\approx0.17$). The direction is right (the content/detail split MCM targets) but the
  magnitude at this weight is modest --- comparable to a well-behaved CVI, without the instability.
\end{itemize}
\textbf{Verdict.} MCM is the \emph{safe} way to nudge the codes toward semantics: it preserves
reconstruction and the mixture, keeps the codebook healthy, and gives a small monotone gain --- but it
is not (yet) the large jump the GAN was. Two levers remain to push it: a \emph{stronger} MSM weight now
that we know it is stable, and the \textbf{generative} sibling --- \emph{masked latent GMM modeling},
which scores the true hidden latent under $p(z\mid\text{ctx})=\sum_k r_k(\text{ctx})\,\mathcal{N}(e_k,
\sigma_k^2 I)$, a conditional version of the code-centered prior. That head is implemented (config
\texttt{dualvae\_msm\_latent\_D.yaml}) and doubles as a \textbf{native generative prior over the GMM
latent} (sample $k\sim r_k$, then $z\sim\mathcal{N}(e_k,\sigma_k^2)$, fill the grid MaskGIT-style,
decode) --- the ``reconstructs at VAE quality, carries a GMM latent, \emph{and} can generate from it''
configuration; its training run is pending.

\part{Verdict}

\section{The overall story}
\label{sec:story}
Read in order, the report is one argument with four load-bearing steps.

\paragraph{1. Claim~A is never in doubt; Claim~B is the whole problem.}
The continuous branch stays active in every run of the design --- both priors, both
$\beta$, with and without the GAN ($100\%$ active units except Run~B's $75\%$, posterior far from
the prior; \S\ref{sec:collapse}) --- and both branches remain load-bearing under ablation
(\S\ref{sec:branches}). The single exception, E+ at $\beta{=}1.0$ ($87.5\%$, one dead channel),
marks the boundary rather than contradicting the rule: it is far outside the useful $\beta$ range
and its reconstruction is broken anyway. Posterior collapse is not a failure mode of this
architecture at any setting one would ship. What \emph{is} in doubt is whether the sum is a
Gaussian mixture.

\paragraph{2. With the standard prior, $\beta$ is the only knob --- and it charges for structure.}
The KL to $\mathcal{N}(0,I)$ pulls the continuous branch toward the \emph{origin}, independently of
the code, so the continuous spread ($\approx1.0$) dwarfs the codeword spacing ($\approx0.1$) and the
codewords are bumps inside one cloud (\S\ref{sec:mixture}).
\begin{itemize}
  \item \emph{Low $\beta$ (Run~A):} the continuous branch runs free, becomes a high-rate
  near-deterministic code and does the reconstruction $\Rightarrow$ \textbf{best fidelity}
  (PSNR $29.4$) but \textbf{no mixture} (separability $0.11$) --- and the codes are largely
  \emph{decorative}: only $\sim$$36\%$ of their apparent contribution is real per-image
  information, the rest a DC/off-manifold artifact of zeroing (\S\ref{sec:codeinfo}).
  \item \emph{Higher $\beta$ (Run~B):} the continuous branch is reined in, the codes spread out and
  carry the image $\Rightarrow$ an \textbf{emerging mixture} (separability $0.11\!\to\!0.68$,
  $\sim$$98\%$ of the code contribution now genuine) bought with reconstruction
  (PSNR $29.4\!\to\!27.3$, rFID $24\!\to\!40$).
\end{itemize}
That is the \textbf{rate/separability trade-off}.

\paragraph{3. The code-centered prior breaks the trade-off --- because it changes the prior's
\emph{shape}, not its strength.}
Re-centering the prior on the codeword, $p(z\mid k)=\mathcal{N}(e_k,\sigma_k^2 I)$, adds the term
$\mu_\Delta^2/2\sigma_k^2$ to the KL, which drives the residual toward the \emph{codeword} instead
of the origin (\S\ref{sec:codeprior}). Every Claim-B metric flips: separability $2.7$/$2.4$, GMM
beats a single Gaussian on BIC \emph{and} held-out likelihood, silhouette positive, dip test
significant, all $256$ codes used --- at \emph{both} $\beta$. Because the mixture no longer needs a
strong KL, one is free to pick low $\beta$ (\textbf{Run~D}) and recover near-best reconstruction
($0.870$ SSIM / $28.9$ PSNR, essentially Run~A) \emph{with} a clean mixture. Run~D is the
best-of-both configuration among the MSE runs.

\paragraph{4. Two controls rule out the alternative explanations --- it really is the prior.}
The mixture could plausibly have come from the codebook, or from adversarial pressure. Neither
survives contact:
\begin{itemize}
  \item \textbf{Not the objective} (\S\ref{sec:twobytwo}, the new A+/B+ runs). Under LPIPS+GAN the
  full $2\times2$ of prior~$\times$~$\beta$ reconstructs \emph{identically} (rFID
  $3.70$--$4.02$, LPIPS $\approx0.043$) --- so the prior is free --- yet the mixture appears in
  exactly the two code-prior cells. With the standard prior the GAN leaves A+ a single blob
  (separability $0.24$) and B+ merely ``emerging'' ($1.19$), with the decisive code-anchored
  likelihood \emph{negative} in both. Worse, without the prior the adversarial objective
  \textbf{kills half the codebook} (A+: $101/256$ codes, $87$ effective components, residual std
  $2.4$); with it, all $256$ stay alive at $226$ effective components.
  \item \textbf{Not the codebook} (\S\ref{sec:vqcontrol}). A plain VQ-VAE sharing D+'s codebook,
  RQ depth, LPIPS and GAN --- but no continuous branch and no code-centered prior --- reconstructs
  well ($5.93$ rFID) yet its codewords merely $k$-means-tile \emph{one} Gaussian: the code-anchored
  held-out likelihood is \textbf{negative} ($-0.97$ vs $+0.16$ for D+). Note this control also shows
  why the geometric proxies alone are not evidence: the VQ-VAE's separability ($3.37$) and
  silhouette ($0.25$) look \emph{better} than D+'s.
  \item \textbf{And not more $\beta$ either} (\S\ref{sec:betaceiling}, Run~E+). The last escape
  route for the standard prior is that $\beta$ was simply too small. Taking it to $1.0$ --- a
  $1000\times$ sweep over A+ --- moves every geometric proxy almost all the way to the code-prior
  values (separability $0.24\to1.19\to2.74$, live codes $101\to190\to256$) and \textbf{still}
  leaves the code-anchored likelihood negative ($-3.12\to-0.18\to-0.08$). The approach is bought
  with rFID $3.70\to4.94$, LPIPS $0.043\to0.066$ and the report's only posterior collapse
  ($87.5\%$ active units). So the rate/separability trade-off of step~2 was never removed by the
  GAN --- only pushed beyond the $\beta$ range we had measured. \emph{$\beta$ sets the latent's
  scale; the prior sets its shape; only shape makes a mixture.}
  \item \textbf{And the failure mode is exactly the one the theory predicts.} Whenever the
  continuous branch's spread outgrows the codeword spacing --- A+ ($2.40$ vs $0.58$), D++ with the
  wavelet branch ($3.15$ vs $2.61$; \S\ref{sec:wavelet}) --- the mixture dies, precisely condition
  \eqref{eq:sepc}. E+ is the same law read the other way: its spacing grows to $4.11$, but so does
  its residual ($1.50$), so the ratio improves without the codewords ever becoming mixture means.
\end{itemize}

\paragraph{What the structure is worth (and what it is not).}
\begin{itemize}
  \item \textbf{It is free.} A matched single-branch ablation at equal capacity and recipe
  (8-ch, LPIPS+GAN, $200$ ep) shows removing the discrete branch barely moves reconstruction
  (D+ $3.97 \to$ VAE $4.08$ rFID --- a tie) while removing the continuous branch clearly hurts
  (D+ $\to$ VQ-VAE $5.93$). So the continuous branch carries reconstruction and the dual model adds
  the Gaussian-mixture latent \textbf{at no reconstruction cost}. The value proposition is a
  structured, code-organized, generatively-usable latent at VAE-level fidelity --- not a pixel win.
  Both single-branch baselines still beat RAE ($7.64$) at equal budget.
  \item \textbf{It is measurably useful, in a narrow way.} The mixture makes the bag-of-codes a
  better linear-probe descriptor (\S\ref{sec:downstream}); the benefit is confined to the
  \emph{discrete} features and to probing --- retrieval is flat everywhere, and the pooled
  \texttt{z\_sum} ties a plain VAE at matched recipe. For downstream semantics the GAN, not the
  GMM, is the bigger lever.
  \item \textbf{It is not object segmentation} (\S\ref{sec:seg}). Painted back onto the image the
  codes track appearance, not objects: every run finds figure/ground on easy images and none
  produces object-level regions. Spatial coherence follows $\beta$ (Run~B is highest) rather than
  the prior; what the prior controls is codebook \emph{usage} ($\approx$$210$ codes/image for
  C+/D+ vs $96$ for A+).
  \item \textbf{Auxiliary semantic objectives: one no-op, one safe small gain.} CVI
  (\S\ref{sec:cvi}) preserved reconstruction and the mixture perfectly but added nothing
  ($0.438\!\to\!0.437$), most likely because its weight is negligible against the sum-reduced
  reconstruction loss. Masked code modeling (\S\ref{sec:mcm}) is the stable alternative: no
  collapse shortcut, reconstruction and mixture intact ($3.97\!\to\!3.85$ rFID), and a small
  monotone gain ($0.438\!\to\!0.441$ bagcodes, $0.357\!\to\!0.361$ \texttt{z\_vq}).
  \item \textbf{A clean negative} (\S\ref{sec:wavelet}): the wavelet detail branch raised rFID
  ($3.9/4.0\!\to\!4.5/4.4$) and weakened or destroyed the mixture. Its gain in
  \texttt{swd\_dualvae} was tied to a band-aware SWD regularizer; the per-code isotropic
  $\sigma_k$ cannot shape three frequency bands.
\end{itemize}

\paragraph{Recommendation.} For a cleanly separated Gaussian-mixture latent, use the
\textbf{code-centered prior} (\texttt{residual\_continuous} + \texttt{component\_prior} +
\texttt{use\_ema\_codebook}) rather than cranking $\beta$, and rather than hoping a perceptual
objective will do it --- it will not. Since the mixture is $\beta$-independent, keep $\beta$
\emph{small} (Run~D/D+) to also get the best reconstruction. Train with LPIPS+GAN: it costs the
mixture nothing at the component level, cuts rFID $\sim$6--8$\times$, and is the single biggest
downstream lever. Further tightening of components is available by lowering
\texttt{sigma2\_floor}. Report \textbf{FID/LPIPS}, not PSNR/SSIM, for the perceptual models, and
judge Claim~B on the \textbf{code-anchored held-out likelihood} --- separability and silhouette
read high even for a plain VQ-VAE that has no mixture at all. When quantifying a branch's
contribution, prefer the \emph{replace-with-mean} intervention over zeroing, to avoid crediting the
branch for a DC/off-manifold artifact. The open thread is the \emph{generative} use of the mixture:
masked latent GMM modeling gives a native prior over this latent (sample $k\sim r_k$, then
$z\sim\mathcal{N}(e_k,\sigma_k^2)$, fill MaskGIT-style, decode).

\vspace{0.5em}
{\sloppy\small
\noindent\textit{All figures and per-run JSON reports (\texttt{latent\_analysis\_*.json},
\texttt{ablation\_report\_*.json}, \texttt{branch\_contribution\_*.json}) are in the
checkpoint folders listed in Table~\ref{tab:ckpts}; downstream results in
\texttt{reports/downstream\_gan/} and the segmentation summary in
\texttt{reports/code\_cluster\_seg\_all\_runs.json}. Reproduce with}
\texttt{tools/dualvae\_latent\_analysis.py}, \texttt{tools/feature\_pca\_images.py},
\texttt{tools/reconstruction\_ablation.py}, \texttt{tools/branch\_contribution.py},
\texttt{tools/linear\_probe\_retrieval.py}, \texttt{tools/code\_cluster\_segmentation.py}.
}

\end{document}
